% The WOP user manual. Build it with `python tools/make_manual.py` (see tools/TOOLS.md), or with pdflatex in this folder.
\documentclass[a4paper,11pt,twoside,openany,parskip=half]{scrbook}
% Packages, colours and the small commands of the WOP manual.

\usepackage[T1]{fontenc}
\usepackage[utf8]{inputenc}
\usepackage[british]{babel}
\usepackage{lmodern}
\usepackage{microtype}
\setlength{\emergencystretch}{3em} % a line may be a little loose instead of too wide (no overfull box)
\usepackage{graphicx}
\usepackage[export]{adjustbox}
\usepackage{xcolor}
\usepackage{booktabs}
\usepackage{tabularx}
\newcolumntype{Y}{>{\raggedright\arraybackslash}X}  % a text column that is not justified
\usepackage{enumitem}
\usepackage{caption}
\usepackage[most]{tcolorbox}
\usepackage{scrlayer-scrpage}

% the version and the date come from tools/make_manual.py (build/version.tex)
\InputIfFileExists{build/version.tex}{}{\newcommand{\wopversion}{unknown}\newcommand{\wopdate}{\today}}

\usepackage[hidelinks, pdftitle={WOP user manual}, pdfsubject={Workshop Optimised Programming for CNC lathes}]{hyperref}

\definecolor{wopblue}{RGB}{40,90,170}
\setkomafont{disposition}{\color{wopblue}\sffamily\bfseries}
\setkomafont{pageheadfoot}{\sffamily\small}
\setkomafont{subject}{\color{wopblue}\sffamily\Large\bfseries}
\setkomafont{subtitle}{\color{wopblue}\sffamily\Large}
\setkomafont{title}{\color{wopblue}\sffamily\bfseries\fontsize{48}{52}\selectfont}
\setkomafont{titlehead}{\sffamily\small\color{black!60}}
\renewcommand*{\chapterheadstartvskip}{\vspace*{0.5\baselineskip}}  % a chapter starts near the top of its page
\renewcommand*{\chapterheadendvskip}{\vspace{1.2\baselineskip}}
\raggedbottom
\setcounter{tocdepth}{2}      % chapters, sections and subsections in the contents
\DeclareTOCStyleEntry[numwidth=3em]{tocline}{section}              % room for numbers like 12.10.
\DeclareTOCStyleEntry[indent=4.5em,numwidth=3.8em]{tocline}{subsection}
\setcounter{secnumdepth}{3}
\captionsetup{font=small,labelfont={sf,bf}}
\setlist{itemsep=0.2em, topsep=0.4em}

\pagestyle{scrheadings}
\clearpairofpagestyles
\ihead{\textsf{WOP user manual}}
\ohead{\headmark}
\cfoot{\pagemark}
\cfoot*{\pagemark}
\automark[section]{chapter}

% --- how things on the screen are written ---------------------------------------------------------------------------------
\newcommand{\menu}[1]{\textsf{\textbf{#1}}}                     % a menu or a menu entry
\newcommand{\menupath}[2]{\menu{#1}\,$\rightarrow$\,\menu{#2}}  % File -> Open...
\newcommand{\button}[1]{\textsf{[#1]}}                          % a button
\newcommand{\field}[1]{\textsf{#1}}                             % the name of a field or a check box
\newcommand{\key}[1]{\fbox{\small\textsf{\strut#1}}}            % a key on the keyboard
\newcommand{\file}[1]{\texttt{\detokenize{#1}}}                 % a file name or a text to type

% --- boxes ---------------------------------------------------------------------------------------------------------------
\tcbset{boxstyle/.style={breakable, boxrule=0.4pt, left=2mm, right=2mm, top=1mm, bottom=1mm, fonttitle=\sffamily\bfseries\small, before skip=0.8em, after skip=0.8em}}
\newtcolorbox{notebox}{boxstyle, colback=blue!4, colframe=blue!50!black, title=Note}
\newtcolorbox{tipbox}{boxstyle, colback=green!5, colframe=green!40!black, title=Tip}
\newtcolorbox{warningbox}{boxstyle, colback=red!5, colframe=red!60!black, title=Warning}
\newtcolorbox{draftbox}{boxstyle, colback=yellow!12, colframe=orange!70!black, title=Draft: not written yet}

% a picture of the program: \shot[width as part of the text width]{file in images/}{caption}
\newcommand{\shot}[3][0.9]{%
  \begin{figure}[htbp]\centering
  \fcolorbox{black!30}{white}{\includegraphics[max width=#1\textwidth, max height=0.62\textheight]{images/#2}}
  \caption{#3}\label{fig:#2}
  \end{figure}}


\begin{document}

\frontmatter
\titlehead{\centering\small Draft\\This manual is machine generated.}
\subject{User manual}
\title{WOP}
\subtitle{Workshop Optimised Programming\\for CNC lathes}
\author{}
\date{Version \wopversion\\\wopdate}
\maketitle

\chapter*{About this manual}
\addcontentsline{toc}{chapter}{About this manual}

This manual is for the people who use WOP in the workshop: the ones who have a part to turn and want a G-code program for the lathe. It explains what each menu, window and button does and how to get a program from the screen to the machine. It does not explain how WOP is built inside; that is in the documents for developers.

\section*{How to read it}
\begin{itemize}
  \item \textbf{Part I} tells you what WOP is, how to start it, and how the main window is laid out.
  \item \textbf{Part II} goes through the menus, one chapter for each menu, in the order of the menu bar.
  \item \textbf{Part III} follows the work: the project, the blocks, contours, tools, workholding, machines, checking and export.
  \item \textbf{Part IV} is for looking things up: keyboard shortcuts, files and folders, a glossary, and what to do when something goes wrong.
\end{itemize}

\section*{How things are written}
\begin{center}
\begin{tabularx}{0.9\textwidth}{@{}l Y@{}}
\toprule
\textbf{Looks like} & \textbf{Means} \\
\midrule
\menupath{File}{Open...} & A menu and an entry in it. \\
\button{Close} & A button in a window. \\
\field{Outer diameter} & The name of a field, a list or a check box. \\
\key{Ctrl}+\key{S} & Keys on the keyboard, pressed together. \\
\file{project.wop.json} & A file name, or text you type. \\
\bottomrule
\end{tabularx}
\end{center}

\begin{notebox}
A note adds something that is good to know.
\end{notebox}
\begin{tipbox}
A tip makes the work easier.
\end{tipbox}
\begin{warningbox}
A warning is about something that can cost you work, a tool or a part.
\end{warningbox}

\section*{Status of this draft}
This manual is machine generated.

WOP is still being developed, and so is this manual. Every chapter is written, but it was written from the program and has not yet been checked by many readers. Where something is not confirmed on a real machine, the text says so. If you find a mistake or something that is missing, please tell the developer. Pictures are taken by a tool from the program itself, so they follow the program when it changes.

\tableofcontents

\mainmatter

\part{Getting started}
\chapter{Introduction}\label{ch:introduction}

\section{What WOP does}
WOP (Workshop Optimised Programming) makes G-code for a CNC lathe. You do not write G-code. You describe the job: the raw part, the shape you want, the tools and what to do in which order. WOP works out the tool paths, shows them in a picture, checks them, and writes the G-code for your controller.

WOP is made for lathes with two axes (X and Z) that run on LinuxCNC. Other controllers are possible through the postprocessors (Chapter~\ref{ch:posts}), but only LinuxCNC is the target.

\section{How a program is made}
A program is a list of \emph{blocks}. Each block is one job on the part: facing, turning, drilling, threading, grooving and so on. The blocks are done from top to bottom. A block works on the material that the blocks above it left over.

You work in this order:
\begin{enumerate}
  \item Set the machine, the raw part and the workholding in the \emph{Setup} (Chapter~\ref{ch:project}).
  \item Add tools (Chapter~\ref{ch:tools}).
  \item Add blocks and fill in their forms (Chapter~\ref{ch:blocks}). A block that follows a shape needs a \emph{contour}, which you draw in the sketch editor (Chapter~\ref{ch:contours}).
  \item Look at the picture, step through the moves and read the messages (Chapter~\ref{ch:checking}).
  \item Export the G-code (Chapter~\ref{ch:file}, Section~\ref{sec:file:export}).
\end{enumerate}
You can go back at any step. WOP plans the whole program again after every change.

\section{The words used in this manual}
\subsection{Project}
A project is one job in one file (\file{.wop.json}). It holds the Setup, the blocks, the contours, and the tools and workholding of the job.

\subsection{Block}
A block is one step of the program. The blocks are listed in the flow list on the left of the window. A block has a form on the right with its values.

\subsection{Contour}
A contour is a shape drawn in the sketch editor, a line of straight lines and arcs. A block can follow a contour. A contour is listed under the heading \emph{Contours} in the flow list. It is not a step of the program, and it can be copied.

\subsection{Tool and workholding}
A \emph{tool} is a cutting tool with its real shape. \emph{Workholding} is the chuck body and the jaws that hold the part. WOP uses their shapes to find collisions.

\subsection{Machine and postprocessor}
A \emph{machine} holds what belongs to your lathe: the travel of the axes, the spindle speed and power, the tool change position and which postprocessor to use. A \emph{postprocessor} (post) is a text file that says how the G-code is written for a controller.

\section{Coordinates and units}
WOP uses millimetres.

\subsection{X as radius or diameter}
X is the distance from the centre line of the part. In the project you choose whether sizes in X are typed and shown as a \emph{radius} or as a \emph{diameter} (Section~\ref{sec:project:xmode}). The labels of the fields say which. The G-code has its own setting in the postprocessor.

\subsection{Z and the picture}
Z runs along the part. The picture shows Z to the right and X up; the chuck is on the left. The face of the part is at $Z=0$ and the part lies towards minus Z.

\section{Where WOP is limited}
\begin{itemize}
  \item WOP has two axes. There is no C axis, no milling and no live tools.
  \item The G-code is written as drawn. WOP does not shorten a path because a tool might not reach; the check reports it and you decide (Chapter~\ref{ch:checking}).
  \item Nothing is tried on every machine. Parts that are not confirmed on a real machine are marked in this manual. Check what a postprocessor writes before you run it; any G-code can crash a machine.
  \item The checks are estimates. The power check is an approximation, and the check of the feed speed is not a full simulation of the controller.
\end{itemize}

\chapter{Installing and starting}\label{ch:installing}

\section{What you need}
\begin{itemize}
  \item A Windows PC (or Linux, see below). A screen of 1024 by 768 pixels is enough.
  \item For the 3D view: a graphics card that can do OpenGL 2. Almost every PC from 2009 on can. Without it, the 3D button is grey and the 2D picture does everything else.
  \item To run the G-code: LinuxCNC on the lathe, and a way to carry the file there.
\end{itemize}

\section{Installing}
WOP for Windows comes as a zip file. Unpack it anywhere on your PC, for example to \texttt{C:\textbackslash WOP}. Nothing is installed. Start only the \file{WOP.exe} that is in the unpacked folder; do not move the file out of it.

The program is not signed. The first time, Windows (SmartScreen) or a virus scanner may warn you. If you got the zip from the place named in the release notes, you can allow it.

\begin{notebox}
On Linux, WOP is installed with a script, and the rest of this manual applies as it is. Developers can also run WOP from the source code with Python. Both are described in the file \file{README.md}.
\end{notebox}

\section{Starting WOP}
Double click \file{WOP.exe}. A window opens with an empty project. Use \menupath{File}{Open...} to open a project file. The example projects in the folder \file{examples} show how finished projects look. They carry the machine of their author (the tool change position and the limits of his lathe): choose your own machine before you export G-code from one of them.

The first time, look at \menupath{File}{Machines...} and make a machine for your lathe (Chapter~\ref{ch:machines}). Then choose it in the Setup.

\section{Where WOP keeps its files}
\subsection{The user folder}
Your settings, machines, libraries and postprocessors are not in the project. They are in the \emph{user folder}. WOP uses one of two places:
\begin{itemize}
  \item If there is a folder \file{user_data} next to \file{WOP.exe}, WOP uses it. Then you can copy the whole program folder to another PC and take everything with you.
  \item Otherwise it uses \texttt{\%APPDATA\%\textbackslash WOP}, in the folder of your Windows user.
\end{itemize}
To change from one to the other, make or remove the folder \file{user_data}. Files are not moved by themselves.

\subsection{Libraries and settings}
The user folder holds the settings (\file{settings.json}), your machines, the libraries of tools, workholding and materials, your own postprocessors, and two log files. Chapter~\ref{ch:app_files} lists them.

\chapter{The main window}\label{ch:window}

WOP has one window. You build the program in it, you see the result in it, and you read what the program found in it. This chapter shows where everything is. The later chapters use the names you learn here.

\shot{window_parts}{The main window. The numbers are explained below: (1) menu bar, (2) flow list, (3) buttons of the flow list, (4) picture controls, (5) picture, (6) messages, (7) form, (8) handles.}

The window has three columns. On the left is the \emph{flow list}: the steps of your program in the order the machine does them. In the middle is the \emph{picture} of the part and the tool paths, and below it the \emph{messages}. On the right is the \emph{form} of the step you chose in the flow list. When you choose another step, the form changes. When you change a value in the form, the picture follows.

The title of the window shows the name of the project file, the name and version of WOP (for example \emph{WOP 0.1.0}), and a star (\file{*}) when there are changes that are not saved. A new project that has no file yet is called \emph{unsaved project}.

\section{The menu bar}\label{sec:window:menus}

The menu bar (1) has five menus.

\begin{center}
\begin{tabularx}{\textwidth}{@{}l Y l@{}}
\toprule
\textbf{Menu} & \textbf{What it is for} & \textbf{Chapter} \\
\midrule
\menu{File} & Projects, G-code, machines, postprocessors, preferences. & \ref{ch:file} \\
\menu{Edit} & Undo and redo. & \ref{ch:edit} \\
\menu{Tools} & The tools of the project and the tool libraries. & \ref{ch:tools_menu} \\
\menu{Workholding} & Jaws and chuck bodies and their libraries. & \ref{ch:workholding_menu} \\
\menu{Materials} & The material table and its libraries. & \ref{ch:materials_menu} \\
\bottomrule
\end{tabularx}
\end{center}

\section{The flow list}\label{sec:window:flow}

The flow list (2) is the program. It is read from the top down, and the machine does the steps in that order.

\subsection{Setup, blocks and contours}
The first entry is always \menu{Setup}. It holds what is true for the whole project: the machine, the raw part, the safe distances and the workholding (Chapter~\ref{ch:project}). It cannot be moved or deleted.

Below it come the \emph{blocks}. A block is one job for one tool, for example facing, or turning along a contour (Chapter~\ref{ch:blocks}). Each block shows

\begin{itemize}
\item a coloured bar at its left edge. The same colour is used for its tool path in the picture, so you can see which path belongs to which block;
\item its number and its type, for example \textsf{3: turning};
\item the name you gave it, if you gave one;
\item a note, if the program has something to say about it. A warning mark and the text \emph{Collision, see messages} mean that the check found a collision in this block.
\end{itemize}

Under the blocks is the heading \menu{Contours}. A \emph{contour} is a shape you drew, for example the finished outline of a part (Chapter~\ref{ch:contours}). A contour is not a step of the program. A block that works along a contour names the contour it uses. Each contour shows how many blocks use it.

\subsection{Choosing an entry}
Click an entry to choose it. The picture and the form show that entry. A click on a contour opens its form at once. A click on Setup or on a block shows its form if the right column is open; a double click opens the right column if it was folded away (Section~\ref{sec:window:handles}).

\subsection{Tick boxes for export}
Each block has a tick box. The box above the list, \field{Tick all blocks for export}, ticks or unticks all of them; it shows a dash when only some are ticked. The ticks matter when you use \button{Export selected...}: only the ticked blocks are written. See Section~\ref{sec:file:selected}. A normal \menupath{File}{Export G-code...} writes all blocks, whatever the ticks say.

\subsection{Adding, moving, deleting and copying}\label{sec:window:buttons}
Five buttons sit under the flow list (3).

\begin{description}
\item[\button{Add block}] opens a list of the block types and, at the end, \menu{Contour}. A new block is put after the block that is chosen, or at the end of the list if none is chosen. \menu{Contour} makes a new contour and opens the sketch editor for it at once (Chapter~\ref{ch:contours}).
\item[\button{Delete}] deletes the chosen block or contour. If other blocks use a contour, WOP asks whether to delete those blocks too or to keep them without a contour.
\item[\button{Up} and \button{Down}] move the chosen block one place up or down. The order of the blocks is the order of the program, and the material that one block removes is gone for the next one.
\item[\button{Export selected...}] writes the ticked blocks (Section~\ref{sec:file:selected}).
\end{description}

A right click on a block or a contour opens a small menu. \menu{Copy} makes a copy of it. For a block there are also two entries that show the operations of the block as plain text (\emph{neutral code}): the moves before they are turned into the language of one controller. They are meant for people who want to see exactly what a block does.

\begin{tipbox}
\menupath{Edit}{Undo} takes back the last change, also a deleted block or a move. If a block does something you did not expect, look at it in the picture before you change it: Section~\ref{sec:window:stepping} shows how to step through the moves.
\end{tipbox}

\section{The picture}\label{sec:window:picture}

The picture (5) shows the raw part, the finished shape, the chuck and the jaws, and the paths of the tool. It is made again after every change.

The controls above it (4) are, from the left: \button{2D} and \button{3D}, \field{Up to selected block}, \field{Step through} and \button{Reset view}. At the right of the picture is the check box \field{Legend}.

\subsection{2D and 3D}
\button{2D} shows a flat section of the part: the radius (X) goes up, the length (Z) goes to the right, and the chuck is on the left. Use the mouse wheel to zoom. Drag with the left or the middle mouse button to move the picture. A double click, or \button{Reset view}, shows the picture as it was at the start. When you zoom in, the lines keep their width, so close paths stay easy to follow.

\button{3D} shows the part as a solid, with the chuck and the jaws, the tool, and the path on the cut face. The shortcut is \key{Ctrl}+\key{3}. The mouse works as in 2D: wheel to zoom, drag to move, double click to start again. The 3D view needs OpenGL. If your computer cannot draw it, the \button{3D} button is grey and its hover text says why.

\subsection{What the picture shows}
Solid lines are cutting moves. Dashed lines are rapid moves. The colours are those of the blocks in the flow list. The raw part is grey and the material that is left after the blocks is brown. The finished shape of the chosen block is drawn dashed in pink. A red cross marks a place where the check found a collision; the messages say what it is.

The 2D picture is fitted to the raw part and the cutting moves, so a small part fills the picture. The tool change position and the safe position are not inside the frame: a rapid move towards them runs out of the picture.

\field{Up to selected block} hides all blocks after the chosen one, so you see the part as it is at that moment of the program.

\subsection{Step through}\label{sec:window:stepping}
\field{Step through} lets you follow the machine move by move. The tool is drawn at the place of the move, and the material that is left at that moment is shown. Four buttons and a slider appear:

\begin{center}
\begin{tabularx}{\textwidth}{@{}l l Y@{}}
\toprule
\textbf{Button} & \textbf{Key} & \textbf{What it does} \\
\midrule
\button{|<} & \key{Home} & Goes to the start of the chosen block. \\
\button{<} & \key{Left} & One move back. \\
\button{>} & \key{Right} & One move forward. \\
\button{>|} & \key{End} & Goes to the end of the chosen block. \\
\bottomrule
\end{tabularx}
\end{center}

The text next to the slider says which block and which move you are at, for example \emph{Block 3, move 12 of 40}. If the tool collides with something in this move, the word \textsf{COLLISION} is added; the messages box has the details.

\begin{notebox}
Stepping through is the best way to find out what a block does. Use it every time you try a block for the first time, before you export.
\end{notebox}

\subsection{Legend and colours}
The legend lists the blocks with their colours and explains the line styles. You can hide it with the check box \field{Legend}; WOP remembers your choice. The colours follow the colours of Windows: with a dark theme the picture is dark.

\section{The form on the right}\label{sec:window:form}

The form (7) shows the values of the chosen entry: Setup, a block, a contour, or, in the tool and workholding views, an item of a library. The title above the form says what it is, for example \emph{Block 3: turning}.

\begin{itemize}
\item A field with a unit shows the unit after the number (mm, mm/rev, deg). You can type the number or use the arrow keys. Both a point and a comma work as the decimal point.
\item Choices are lists. A check box turns something on or off. A field that says \emph{not set} is optional: leave it empty to use the value of the block.
\item A \button{Dia.} or \button{Rad.} button next to a field changes how a size in X is shown: as a diameter or as a radius (Section~\ref{sec:project:xmode}).
\item A field that is grey and cannot be changed shows a value the program works out. Example: the Start Z of a turning block.
\end{itemize}

\subsection{Values you have not touched}
When a new block is made, its fields get preset values: the defaults of the Preferences, or the values that belong to the tool you chose. A preset value is shown in a lighter colour. As soon as you type a value in a field, it gets the normal colour, and WOP no longer changes it by itself. If you change the tool, the values of the fields you did not touch follow the new tool; the values you typed stay.

A right click on a field opens a menu with \menu{Reset to default}. The field gets its preset value again and the light colour.

\subsection{The line under the form}
Under the form of the Setup, one line says how long the raw part is and how much of it sticks out of the chuck.

\section{The messages}\label{sec:window:messages}

The box under the picture (6) is where the program tells you things. A line starts with a word that says how serious it is.

\begin{center}
\begin{tabularx}{\textwidth}{@{}l Y@{}}
\toprule
\textbf{Starts with} & \textbf{Meaning} \\
\midrule
\textsf{note:} & Something you should know, for example that a part of the material is not reached by the tool. Nothing is wrong. \\
\textsf{COLLISION:} & The check found a collision: the tool, its holder or the jaws touch material or the part where they must not. The line says where and gives a way out. \\
\textsf{PROBLEM:} & The program cannot be planned. The text says why. \\
\textsf{WORKING:} & The program is being planned or checked. The picture and the notes are those of the program before your last change until the work is done. \\
\textsf{NOTICE:} & A short message about what you just did, for example that a value was refused. \\
\bottomrule
\end{tabularx}
\end{center}

Drag the border between the picture and the messages (the three dots) up or down to give more room to the one or the other. If there is nothing to say, the box shows \emph{No problems found.} The status bar at the bottom of the window shows the same short messages for a few seconds, and says \emph{Planning ...} or \emph{Checking collisions and power ...} while the work goes on. You can keep working while that goes on; the window stays usable.

\section{Panels and handles}\label{sec:window:handles}

The two side columns can be made wider or narrower, or folded away to give the picture more room. The handles between the columns (8) show three dots. Drag a handle to change the width. Click it once, without dragging, to fold the column away; a folded column leaves a handle with an arrow, and a click on it opens the column again. WOP remembers the sizes when you close it.

WOP is made to work on a screen of 1024 by 768 pixels, so it also fits the small screen of a machine.


\part{The menus}
\chapter{The File menu}\label{ch:file}

The \menu{File} menu is where a project starts and ends. It opens and saves projects, writes the G-code, and gives you the dialogs for machines, postprocessors and preferences. This chapter goes through the entries from the top down.

\shot[0.32]{menu_file}{The \menu{File} menu.}

\begin{center}
\begin{tabularx}{\textwidth}{@{}l l Y@{}}
\toprule
\textbf{Entry} & \textbf{Keys} & \textbf{What it does} \\
\midrule
\menu{New} & \key{Ctrl}+\key{N} & Starts a new project. \\
\menu{Open...} & \key{Ctrl}+\key{O} & Opens a project from a file. \\
\menu{Save} & \key{Ctrl}+\key{S} & Saves the project. \\
\menu{Save as...} & \key{Ctrl}+\key{Shift}+\key{S} & Saves the project under a new name. \\
\menu{Export G-code...} & \key{Ctrl}+\key{E} & Writes the G-code of the whole program. \\
\menu{Export selected blocks...} & & Writes the G-code of the blocks you ticked. \\
\menu{Machines...} & & Opens the list of machines. \\
\menu{Postprocessors...} & & Opens the postprocessor editor. \\
\menu{Preferences...} & & Opens the settings that apply to every project. \\
\menu{Quit} & \key{Ctrl}+\key{Q} & Closes WOP. \\
\bottomrule
\end{tabularx}
\end{center}

\section{Projects}\label{sec:file:projects}

A \emph{project} is one job: the part, the program for it and what it needs. WOP works on one project at a time. A project is kept in one file whose name ends in \file{.wop.json}. The file holds the project settings, the raw part, the blocks, the contours, and the tools and workholding the project uses. The libraries (all your tools, workholding and materials) are kept apart from the project, because they belong to you and not to one job.

\subsection{New}
\menupath{File}{New} starts an empty project. It begins with the values you set in the Preferences (Section~\ref{sec:file:preferences}): the material, the size of the raw part and the safe distances. If the project you are working on has unsaved changes, WOP asks first (Section~\ref{sec:file:questions}).

\subsection{Open}
\menupath{File}{Open...} shows the file dialog of Windows. It lists project files, that is files that end in \file{.wop.json} or \file{.json}. The dialog starts in the folder of the last file you opened or saved, so you do not have to look for your folder every time. If a file cannot be read, WOP says why and keeps the project you had.

\subsection{Save and Save as}
\menupath{File}{Save} writes the project to its file. A project that has no file yet asks for a name first, the same as \menupath{File}{Save as...}. \menu{Save as} always asks for a name and suggests \file{project.wop.json} in the folder of your last file. After \menu{Save as}, the new file is the file of the project.

\begin{tipbox}
Save often. \menupath{Edit}{Undo} can take back your last changes, but it cannot bring back a project that was never saved.
\end{tipbox}

\subsection{Questions before something is lost}\label{sec:file:questions}
\menu{New}, \menu{Open...} and \menu{Quit} all end the project you are working on. WOP makes sure you do not lose anything:
\begin{enumerate}
  \item If the project has unsaved changes, it asks \emph{``Save the changes first?''} The answers are \button{Save}, \button{Discard} and \button{Cancel}. \button{Cancel} stops what you started, and nothing is changed.
  \item If tools of the project differ from your tool library, or are not in it yet, it asks whether to update the library. It lists the tools and offers \button{Update the library}, \button{Leave the library as it is} and \button{Cancel}. Choose \button{Update the library} when you changed a tool on purpose and want the better version in all your future projects.
\end{enumerate}

\section{Exporting G-code}\label{sec:file:export}

A G-code program only helps when it is written the way you want it and can be trusted. So before WOP writes a file, it checks the program and asks when it found something. WOP follows one rule here: \emph{the G-code is written as drawn}. It never shortens or skips a path because the tool might not reach; if the picture shows a problem, the G-code will do it too. That is why you are asked.

Open the menu with a click on \menu{File} and choose the entry (Figure~\ref{fig:file_menu_in_window}).

\shot[0.55]{file_menu_in_window}{Choosing \menupath{File}{Export G-code...}. The red frame marks the menu, the pointer the entry.}

\subsection{Export G-code}
\menupath{File}{Export G-code...} writes the whole program. It goes in these steps:

\begin{enumerate}
  \item \textbf{The checks.} WOP looks for collisions and works out the power the cuts need. The checks run by themselves after every change, and the export makes sure they are done. Then it looks at what they found:
  \begin{itemize}
    \item Nothing was found: WOP goes straight on to the next step.
    \item Problems were found: a window lists them (Figure~\ref{fig:export_found}). \button{Export anyway} writes the G-code as it is drawn. \button{Cancel} goes back to the program so you can correct it.
  \end{itemize}
  \item \textbf{What the controller cannot do.} The postprocessor of your machine may lack something a block needs, for example a dwell. Then a window titled \emph{``The controller cannot do everything''} lists each case with a way out and a check box. Tick the ways out you accept. A block that has no accepted way out is left out of the G-code, and the list says so.
  \item \textbf{The file.} A file dialog suggests a name: the title of the project if you gave one in the Setup, else the name of the project file, with the extension of the postprocessor (\file{.ngc} for LinuxCNC). The dialog starts in the folder of your last export. The file is a plain text file.
\end{enumerate}

\shot[0.62]{export_found}{The checks found problems. The G-code is written as drawn, so the tool would really do this. Blocks that are left out are listed in the same window.}

When the file is written, the message line at the bottom of the window says which file was written and with which postprocessor. If blocks were left out, the message adds \emph{``INCOMPLETE''}, and the program itself starts with the comment \emph{``INCOMPLETE - DO NOT RUN''} followed by the reason. Read that line before you take the file to the machine.

\begin{warningbox}
WOP checks collisions and power with the models it has of your tools, workholding and material. It cannot know everything about your machine. Look at the picture, look at the G-code, and make an air cut before the first real cut of a new program.
\end{warningbox}

\subsection{Export selected blocks}\label{sec:file:selected}
Sometimes you want the G-code of a few blocks only, for example to try one cut or to run the facing first. In the flow list on the left, every block has a tick box. Tick the blocks you want, or use \field{Tick all blocks for export} above the list. Then choose \menupath{File}{Export selected blocks...}. If no block is ticked, WOP only says \emph{``Tick at least one block first.''}

\shot{export_selected}{Export selected blocks. On the left the picture and the notes, on the right the G-code that will be written.}

The window shows what you will get before you save it (Figure~\ref{fig:export_selected}):
\begin{itemize}
  \item At the top you say how the blocks you did \emph{not} tick count: \button{done} means their material is gone (they were run before), \button{not done} means their material is still there. This decides how the ticked blocks are drawn and what is checked, and the G-code starts with a comment that tells which blocks it holds.
  \item The picture on the left shows the paths of the ticked blocks on the part. Under it, the notes list problems and remarks.
  \item The G-code on the right is the text that will be written.
  \item \button{Save G-code...} asks for the file name as described above. \button{Close} closes the window without saving.
\end{itemize}

\section{Machines, postprocessors and settings}

\subsection{Machines}
\menupath{File}{Machines...} opens the list of your machines (Figure~\ref{fig:machines}). A machine holds what belongs to one lathe: its postprocessor, the usual position for changing tools, and its limits. The project uses one machine, which you choose in the Setup of the project. The project keeps its own copy of it, so the project still works on a computer that does not have the machine in its list.

\begin{center}
\begin{tabularx}{\textwidth}{@{}l Y@{}}
\toprule
\textbf{Button} & \textbf{What it does} \\
\midrule
\button{New} & Adds a new machine to the list. \\
\button{Import...} & Reads a machine from a file into the list. \\
\button{Export...} & Writes the chosen machine to a file, for example to give it to someone. \\
\button{Delete} & Takes the chosen machine out of the list. A project that uses it keeps its own copy. \\
\button{Make default} & New projects start with this machine. \\
\button{Use in project} & The open project gets a copy of this machine. \\
\button{Add project's machine} & Puts the copy of the machine that is kept in the open project into your list. This is for a project that you opened on another computer, where the machine is not in the list. \\
\button{Close} & Closes the dialog. \\
\bottomrule
\end{tabularx}
\end{center}

The settings of a machine are described in Chapter~\ref{ch:machines}.

\subsection{Postprocessors}
A \emph{postprocessor} (``post'') turns the program into the G-code of one kind of controller. \menupath{File}{Postprocessors...} opens the editor for posts (Figure~\ref{fig:post_editor}). WOP comes with the post for LinuxCNC. Posts are text files, and a new controller is a new file, not a new program. In the folder \file{examples/posts} there are reference posts for other controllers that you can import and change (Chapter~\ref{ch:posts}).

At the top of the window you choose the post and manage the files with \button{New}, \button{Save}, \button{Save as copy...}, \button{Import...}, \button{Delete} and \button{Folder}. The posts that come with WOP cannot be changed. Use \button{Save as copy...} to make your own version, and change that.

\begin{warningbox}
Every post, the ones that come with WOP too, can crash a machine. Check what a post writes before you run it on a machine. The editor says the same at the bottom of its window.
\end{warningbox}

The pages and settings of a post are described in Chapter~\ref{ch:posts}.

\subsection{Preferences}\label{sec:file:preferences}
\menupath{File}{Preferences...} lists the values that WOP uses when you start something new (Figure~\ref{fig:preferences}): the material and size of the raw part of a new project, the safe distances, the feed, depth of cut and allowance of a new block, and values used for the tool paths. For each value the dialog shows what it means and, in grey below it, the name of the value and what is allowed. A value that has a list of choices, such as the material, is a dropdown: click it, then click the entry you want (Figure~\ref{fig:preferences_dropdown}).

\shot[0.95]{preferences_dropdown}{A dropdown in the preferences. Click the value (1) and the list opens; click the entry you want.}

\shot[0.8]{preferences}{The preferences. The sections are \emph{New project}, \emph{New blocks}, \emph{Tool paths} and more below.}

\begin{itemize}
  \item A value you have not changed is shown in a lighter colour. Once you change it, the text becomes normal.
  \item \button{Reset} next to a value puts back the value WOP comes with.
  \item \button{Fixed values...} shows the values that cannot be changed here, such as tables from standards, as text.
  \item Changes are saved at once, but they apply \emph{after you restart WOP}. WOP does not restart by itself, so save your project first.
\end{itemize}

The preferences belong to you and to this computer, not to a project. A project you send to someone else does not carry them.

\section{The sketch editor}
There is no entry for the sketch editor in the \menu{File} menu. It opens from the button \button{Sketch editor...} of a contour, of a tool and of a workholding item, and by itself when you add a contour (Figure~\ref{fig:sketch_editor}). It edits that contour or shape, and \button{OK} stores the result in the project. You can draw lines, arcs and circles, add relations and dimensions, and import a DXF drawing. The buttons have a pictogram and no text. Rest the mouse pointer on a button for a moment and its name appears next to it (Figure~\ref{fig:sketch_hover_text}). How to draw and dimension is explained in Chapter~\ref{ch:contours}.


\section{Quit}
\menupath{File}{Quit} closes WOP. It asks the same two questions as \menu{New} and \menu{Open...} (Section~\ref{sec:file:questions}), so you never lose a change by accident. When WOP closes, it remembers the place and size of the window and the width of the panels, and opens the same way next time. Every dialog (machines, postprocessors, toolboxes, preferences and the others) also opens with the size and place it had when you closed it.

\chapter{The Edit menu}\label{ch:edit}

The \menu{Edit} menu has two entries.

\begin{center}
\begin{tabularx}{\textwidth}{@{}l l Y@{}}
\toprule
\textbf{Entry} & \textbf{Keys} & \textbf{What it does} \\
\midrule
\menu{Undo} & \key{Ctrl}+\key{Z} & Takes back the last change. \\
\menu{Redo} & \key{Ctrl}+\key{Y} or \key{Ctrl}+\key{Shift}+\key{Z} & Does again what Undo took back. \\
\bottomrule
\end{tabularx}
\end{center}

\section{Undo}
\menupath{Edit}{Undo} takes back your last change to the project: a value in a form, a block you added, moved or deleted, a change in the Setup. You can go back several steps.

\section{Redo}
\menupath{Edit}{Redo} puts back what Undo took away. After you make a new change, the steps that were undone cannot be redone any more.

\begin{notebox}
Undo works on the project. It does not bring back a project you did not save, and it does not change the libraries. The sketch editor has its own undo for the drawing (Chapter~\ref{ch:contours}).
\end{notebox}

\chapter{The Tools menu}\label{ch:tools_menu}

The \menu{Tools} menu is the way to your tools. What a tool is, and how its form works, is in Chapter~\ref{ch:tools}.

\begin{center}
\begin{tabularx}{\textwidth}{@{}l Y@{}}
\toprule
\textbf{Entry} & \textbf{What it does} \\
\midrule
\menu{Show tools} & Shows the list of the tools of the project in the middle of the window. \\
\menu{Add tool} & Adds a tool of the type you choose. \\
\menu{Use toolbox} & Puts a toolbox into the project. \\
\menu{Edit toolboxes...} & Opens the window to make and change toolboxes. \\
\menu{Library} & Chooses and manages the tool library. \\
\bottomrule
\end{tabularx}
\end{center}

\section{Show tools}
\menupath{Tools}{Show tools} replaces the picture in the middle by the \emph{Tools view}. It has the list of the tools of the project, a picture of the chosen tool, and the buttons to add, delete and draw tools. The form of the chosen tool is on the right. To get the picture of the program back, choose a block in the flow list.

\section{Add tool}
\menupath{Tools}{Add tool} lists the types of tool: turning tool, turning tool (ISO holder), boring bar, boring bar (ISO), grooving tool, threading tool and drill. A new tool of the chosen type gets the next free number and standard values. The same list is on the button \button{Add tool} of the Tools view.

\section{Use toolbox and Edit toolboxes}
A \emph{toolbox} is a set of tools from the library with the numbers they have in the project, for example the tools that sit in the turret of your lathe. \menupath{Tools}{Use toolbox} lists the toolboxes of the library; the default one is marked \emph{(default)}. If a toolbox would change the tool of a block, WOP lists those blocks and asks before it goes on. A toolbox replaces all tools of the project.

\menupath{Tools}{Edit toolboxes...} opens the window of the toolboxes (Section~\ref{sec:tools:toolboxes}).

\section{Library}\label{sec:tools_menu:library}
The \menu{Library} submenu is the same for tools, workholding and materials. It chooses which library is in use, and it manages the libraries:
\begin{itemize}
  \item The list of libraries. The one in use has a check mark. Click another one to use it.
  \item \menu{New library...} makes an empty library.
  \item \menu{Save as library...} saves the one in use under a new name.
  \item \menu{Import library file...} copies a library file of somebody else into your libraries.
  \item \menu{Open library folder} shows the folder in Windows.
  \item \menu{Rename this library...} and \menu{Delete this library...}. Both are off for the library \emph{Example}, which is always there. When you delete a library, its file is removed and the Example is used instead.
\end{itemize}
A library holds the tools you may want in many projects. A project keeps its own copy of the tools it uses, so changing a library never changes a finished project.

\chapter{The Workholding menu}\label{ch:workholding_menu}

The \menu{Workholding} menu is for the chuck body and the jaws. Chapter~\ref{ch:workholding} explains them.

\begin{center}
\begin{tabularx}{\textwidth}{@{}l Y@{}}
\toprule
\textbf{Entry} & \textbf{What it does} \\
\midrule
\menu{Show workholding} & Shows the list of workholding items in the middle of the window. \\
\menu{Add workholding} & Adds a \emph{chuck body} or \emph{chuck jaws}. \\
\menu{Library} & Chooses and manages the workholding library. \\
\bottomrule
\end{tabularx}
\end{center}

\section{Show workholding}
\menupath{Workholding}{Show workholding} shows the Workholding view in the middle: the items of the library, a picture of the chosen item and the buttons \button{Add workholding}, \button{Delete} and \button{Sketch editor...}. The form of the chosen item is on the right.

The text under the picture reminds you that the jaws and the chuck body the program uses are chosen in the Setup. If none is chosen there, nothing is drawn and the checks use the flat clamping plane of the Setup.

\section{Add workholding}
\menupath{Workholding}{Add workholding} offers \emph{Chuck body} and \emph{Chuck jaws}. The new item has a placeholder shape; you reshape it in the sketch editor.

\section{Library}
The \menu{Library} submenu works as described for tools (Section~\ref{sec:tools_menu:library}). Workholding is kept in a library of its own, so you can have one library for each of your lathes or chucks.

\chapter{The Materials menu}\label{ch:materials_menu}

The material of the raw part decides how much power the cut needs. The \menu{Materials} menu has the table of materials.

\begin{center}
\begin{tabularx}{\textwidth}{@{}l Y@{}}
\toprule
\textbf{Entry} & \textbf{What it does} \\
\midrule
\menu{Edit materials...} & Opens the material table. \\
\menu{Library} & Chooses and manages the material library. \\
\bottomrule
\end{tabularx}
\end{center}

\section{Edit materials}
\menupath{Materials}{Edit materials...} opens the window \emph{Materials}. The list of materials is on the left, with the buttons \button{New} and \button{Delete}. On the right are the \field{Name} and the \field{Specific cutting force} of the chosen material, in N/mm$^2$. The built-in materials are mild steel, stainless steel, aluminium, brass and cast iron.

A change is saved when you leave the field. If the value cannot be saved (an empty name, a name that exists already), the reason stands under the form and the value stays in the field until you correct it.

\begin{itemize}
  \item \button{New} makes a material called \emph{New material}; give it a name and a value.
  \item \button{Delete} removes a material of yours. For a built-in material that you changed, the same button asks whether to bring back its built-in values. If the raw part uses the material, the question says so.
\end{itemize}
A material that you added or changed carries a mark in the list.

\begin{notebox}
The specific cutting force is the force per cut area that the material resists with. The values of the built-in materials are guesses from tables, not measured on your machine. The power check uses them only as an estimate (Section~\ref{sec:checking:power}).
\end{notebox}

\section{Library}
The \menu{Library} submenu works as described for tools (Section~\ref{sec:tools_menu:library}). Use a library of your own if you work with materials that are not in the built-in list.


\part{Making a program}
\chapter{The project and the Setup}\label{ch:project}

A \emph{project} is one job (Section~\ref{sec:file:projects}). Its first entry in the flow list is \menu{Setup}. The Setup holds everything that is true for the whole program and does not belong to one block: which machine, what the raw part is like, where the tool may move safely, and how the part is held. Fill it in first. The blocks after it use these values.

\shot[0.6]{setup_form}{The form of the Setup: project settings, part, safe distances and workholding.}

The form has four groups. The headings in the form are \emph{Project settings}, \emph{Part}, \emph{Safe distances} and \emph{Workholding}.

\section{Project settings}\label{sec:project:settings}

\subsection{Title}
The title of the project. WOP writes it as a comment at the top of the G-code, and it suggests it as the name of the G-code file when you export (Section~\ref{sec:file:export}). If the title is empty, the name of the project file is used.

\subsection{Machine}\label{sec:project:machine}
The machine the program is for. The list shows the machines you made in \menupath{File}{Machines...} (Chapter~\ref{ch:machines}) and the entry \emph{(none)}. The machine decides which postprocessor writes the G-code, where the tool change position is, and what the machine can do: its speeds and its travel. WOP checks the program against these limits.

The button \button{Machine settings...} opens the same dialog as \menupath{File}{Machines...}.

A project always keeps its own copy of the machine, so a project still opens on another PC that does not have this machine in its list. If you open a project that has no machine, WOP gives it the default machine of your list, and says so in the status bar. The project then counts as changed; save it to keep the machine.

\subsection{X mode}\label{sec:project:xmode}
X mode says how you think of X on your lathe: \emph{diameter} or \emph{radius}. It decides three things.

\begin{itemize}
\item How sizes in X are shown and typed in the forms. With \emph{diameter}, a size like the target of a turning block is a diameter, and the label says so.
\item In which mode the G-code is written. LinuxCNC has two modes: \file{G7} (X is a diameter) and \file{G8} (X is a radius).
\item The first value of the \button{Dia.} and \button{Rad.} buttons that sit next to many X fields.
\end{itemize}

The button \button{Dia.} or \button{Rad.} next to a field changes only that one field, and only what is shown. WOP keeps all sizes as radii inside, so nothing changes in the program when you switch the button; you only see the same value the other way.

\begin{notebox}
Typing a diameter where WOP wants a radius is the most common mistake. When a result looks twice as big or twice as small as you meant, look at the label and at the \button{Dia.} / \button{Rad.} button of that field.
\end{notebox}

\subsection{Work offset}
The work offset that the program uses (for LinuxCNC \file{G54} to \file{G59.3}). The list is the work offsets of the postprocessor of the machine (Chapter~\ref{ch:posts}). The G-code selects it in its header. It is the zero point of the part on the machine, which you set up on the lathe itself. A postprocessor that writes no work offset leaves the field grey. If a project has an offset that the post of the machine does not know, the list shows it with the note \emph{not in the post}, and the G-code still uses it.

\section{The raw part}\label{sec:project:part}

\subsection{Material}
The material of the part. WOP takes it from the material table (Chapter~\ref{ch:materials_menu}) and uses it to estimate the cutting power. If the estimate for a block is higher than the machine can give, the check says so (Chapter~\ref{ch:checking}).

\subsection{Bar or tube}
The raw part starts as a plain bar or tube.

\begin{description}
\item[\field{Outer diameter}] The diameter of the bar.
\item[\field{Inner diameter}] The diameter of the hole in a tube. Leave it \file{0} for a solid bar.
\item[\field{Bar stock (unlimited length)}] Ticked: the bar is long and runs through the spindle, so there is material behind the chuck as far as you ever need. Not ticked: the part has a fixed length. Then \field{Length} is the whole length of the part, also the piece inside the chuck, and you can type it. For bar stock, \field{Length} is grey.
\end{description}

The line under the form says in words what you set: \emph{Bar stock: the length is not limited}, or the length, how much sticks out and how much is in the chuck.

\subsection{Drawing the raw part}\label{sec:project:draw}
Not every raw part is a plain bar. A forging, a casting, or a part that was turned in an earlier set-up has steps, a bore or a taper. For these, draw the shape. Click \button{Draw the raw part...}. The sketch editor opens (Chapter~\ref{ch:contours}) with your bar or tube already drawn, so you change a shape that exists and do not start from an empty sheet.

\shot{sketch_raw_part}{The sketch editor opened for the raw part. The green dashed line is the clamping plane. The bar was changed to a stepped shaft.}

Draw the half section of the piece, as for a turning contour: X is the radius and goes up, Z is the length and goes to the right, the spindle axis is the dash-dotted line. The sketch must follow these rules:

\begin{itemize}
\item \textbf{A solid piece} is one open outline. It starts on the spindle axis, goes round the piece, and ends on the spindle axis. The axis closes it.
\item \textbf{A hollow piece} (a tube, a ring, a casting with a bore) is one closed loop that stays above the axis. If a loop touches the axis, WOP treats it as a solid piece.
\item \textbf{The clamping plane.} One point must have the name \emph{Clamping plane}. It marks where the jaws begin; only its Z matters. It must lie inside the length of the piece. Use the button with this name in the section \emph{Mark} of the sketch editor (Section~\ref{sec:sketch:names}). A thin green dashed line across Z shows the plane in the drawing.
\item The outline must not cross or touch itself.
\end{itemize}

Draw the \emph{whole} piece, also the part that is inside the jaws. Only the part outside the jaws is material for the blocks. The text under the drawing says what is still wrong, and the \button{OK} button stays off until the sketch follows the rules.

When you press \button{OK}, the numbers of the Setup follow the drawing: \field{Outer diameter} is the biggest diameter, \field{Inner diameter} is the smallest bore if the piece is hollow, \field{Stick-out} is the distance from the front to the clamping plane, \field{Length} is the whole length, and \field{Front z} is the front of the piece. These fields are grey now, because the drawing is the truth. You can still change \field{Material}. The safe distances follow the drawing, as long as you did not type them yourself (Section~\ref{sec:project:safe}).

\shot[0.6]{setup_raw_part}{The Setup after the raw part was drawn: the numbers of the part are grey, and a button goes back to the plain bar.}

To change the shape later, click \button{Draw the raw part...} again. The button \button{Back to a plain bar or tube} forgets the drawing. The numbers it gave stay in the fields, and you can edit them again.

\begin{notebox}
A drawn raw part with a stepped bore uses the \emph{smallest} bore for \field{Inner diameter}. This number is only used for the check of inside jaws (Section~\ref{sec:project:holding}).
\end{notebox}

\section{Safe distances}\label{sec:project:safe}

These five values say where the tool may move without touching anything. Every block begins and ends at these places, and WOP moves the tool there between blocks.

\begin{description}
\item[\field{Safe position, X radius} and \field{Safe position, Z}] The place where the tool waits between blocks. It must be outside the part and the chuck. Blocks travel at this radius.
\item[\field{Tool change position, X radius} and \field{Tool change position, Z}] The place where the tool is changed. The tool goes here when the next block uses another tool. If your machine has a fixed tool change position in machine coordinates (set in the machine dialog), that position is used and written with \file{G53}.
\item[\field{Clearance (X and Z)}] The distance the tool keeps from the material while it moves inside a block, the same for X and Z.
\end{description}

A new project fills these values in for you from the size of the part: the safe radius is the outer radius plus a margin, and the safe Z is the front of the part plus a margin. The margins are in the Preferences (Section~\ref{sec:file:preferences}). The values \emph{follow the part} as long as you have not typed them: if you make the bar thicker, the safe radius grows with it. As soon as you type one of these values yourself, it stays as you typed it. The tool change position also follows the machine, if the machine has a position of its own.

\begin{warningbox}
Look at the safe position once on the real machine. WOP cannot know where your turret, your tailstock or your steady rest are. If the safe position is not really safe on your lathe, move it here.
\end{warningbox}

\section{Workholding in the Setup}\label{sec:project:holding}

The last group says how the part is held.

\begin{description}
\item[\field{Stick-out (free length)}] How far the part sticks out of the jaws, measured from the front of the part. It places the clamping plane. With a drawn raw part this field is grey; the clamping plane of the drawing places it.
\item[\field{Front z}] The Z of the front face of the raw part. It is usually \file{0}, and the length of the part runs to negative Z.
\item[\field{Workholding}] The jaws or the chuck body from your workholding library (Chapter~\ref{ch:workholding}). \emph{(none)} means that no chuck is drawn and WOP only checks against a flat clamping plane. With jaws, the chuck body that belongs to them is used too.
\end{description}

When you choose jaws, two more fields appear.

\begin{description}
\item[\field{Jaw point on the part}] Which named point of the jaws sits on the part. \emph{(as drawn)} leaves the jaws where they were drawn. A point like \emph{step 1} puts that step of stepped jaws on the part.
\item[\field{Grip inside (tube held from the bore)}] The jaws expand in the bore of a tube and hold it from the inside.
\end{description}

WOP checks the way the part is held and tells you in the messages if something does not fit. For example: the stick-out is longer than the part, a bar is too thick to go through the chuck body, or the jaws of an inside grip reach into the wall of the tube.

\begin{notebox}
WOP does not calculate how far a part can stick out of the jaws without bending or vibrating. That is for you to decide.
\end{notebox}

\chapter{Blocks}\label{ch:blocks}

A program in WOP is a list of \emph{blocks}. A block is one job that is done with one tool: facing the end, turning a shoulder, drilling a hole, cutting a thread. WOP plans the moves of each block from the numbers you give it, checks them, and writes them as G-code. This chapter explains what all blocks have in common and then each type of block.

\section{Adding, moving, deleting and copying blocks}\label{sec:blocks:flow}

Blocks are made and ordered in the flow list (Section~\ref{sec:window:buttons}). \button{Add block} opens the list of types and puts the new block after the one that is chosen. \button{Up} and \button{Down} change the order, \button{Delete} removes a block, and a right click with \menu{Copy} makes a copy.

\emph{The order matters.} The material that one block cuts away is gone for the next one. A roughing block first and a finishing block after it is therefore not two jobs on the same bar: the finishing block starts from what the roughing block left, and its plan uses that. If you move the blocks round, WOP plans them again in the new order.

A new block gets sensible values: the defaults from the Preferences, and for the depth of cut and the allowance the values that belong to the tool you chose. Change what you need. A value you have not touched is shown lighter and follows a change of tool (Section~\ref{sec:window:form}).

\section{What every block has}\label{sec:blocks:common}

Most blocks start with the same fields.

\subsection{Name}
A name of your own. The flow list shows it under the type, so you can tell two turning blocks apart (for example \emph{rough} and \emph{finish}). It is also written as a comment in the G-code.

\subsection{Tool}
The tool of this block, chosen from the tools of the project (Chapter~\ref{ch:tools}). The list shows the tool number, its name and its description. \emph{No tool chosen} leaves the block without a tool. A block needs a tool to be planned. When you change the tool, the depth of cut and the allowance follow the tool as long as you have not typed them.

\subsection{Spindle}
Four fields describe the spindle.

\begin{description}
\item[\field{Spindle: Direction}] \emph{clockwise} or \emph{counter-clockwise}. A new block always starts clockwise. WOP writes the right code for the controller of your machine (the postprocessor does this, Chapter~\ref{ch:posts}).
\item[\field{Spindle: Mode}] \emph{rpm} keeps the number of revolutions fixed. \emph{css} (constant surface speed) keeps the cutting speed fixed and lets the spindle speed change with the diameter.
\item[\field{Spindle: Value}] The number of revolutions per minute for \emph{rpm}, or the cutting speed in metres per minute for \emph{css}.
\item[\field{Spindle: Rpm limit}] Only for \emph{css}: the highest spindle speed that is allowed. While the tool goes to a small diameter, a constant surface speed asks for a very high spindle speed. The limit holds it back. Leave it empty for no limit.
\end{description}

\subsection{Feed}
The feed of a cut, in millimetres per revolution (mm/rev). It is the distance the tool moves for each turn of the spindle. Some blocks have a second field for the finishing pass.

\subsection{Coolant}
\emph{none}, \emph{flood} or \emph{mist}. WOP writes the code of the coolant when it changes from one block to the next (switching it on, to another kind, or off), if the postprocessor of your machine knows these codes.

\subsection{Depth of cut and finish allowance}
Cutting blocks take their material off in several passes. The \field{Depth of cut} is how deep each roughing pass goes, measured in X. The \field{Finish allowance} is the layer of material that roughing leaves on the part for the finishing pass.

\subsection{The finishing pass}
A block with a \field{Finishing pass} check box can cut the last layer in a separate pass that follows the shape exactly. When it is ticked, you can give the pass its own feed and its own spindle speed. An empty field means that the value of the block is used.

\subsection{Cutter compensation}
The turning blocks have a check box \field{Cutter compensation}. When it is ticked, the finishing pass is written for the control to correct for the radius of the tool nose (\file{G41} and \file{G42}), if the postprocessor of the machine can write it. The pass gets a lead-in and a lead-out for it. Without it, the path in the code is already offset for the nose radius by WOP. The LinuxCNC postprocessor that comes with WOP writes it; it was tried on the lathe in an air cut, not yet in material. Not every postprocessor writes it.

\subsection{How WOP plans what you draw}
Four rules hold for all cutting blocks. They are important, so read them once.

\begin{itemize}
\item \textbf{WOP cuts what you draw.} The shape comes from the numbers in the form and from the contour. WOP never shortens or skips it because the tool might not reach. If a block cannot make something, it says so and leaves it as it is.
\item \textbf{Roughing stays off the finished part.} Roughing is planned from the material that is there, the finished shape, the allowance and the depth of cut. It keeps the body of the tool off the part and goes at most 30 percent of the allowance into it. A line that the tool cannot cut all the way is cut as far as it reaches, or not at all. What stays is left for another tool, and a note says so.
\item \textbf{Finishing follows the shape.} The finishing pass runs along the whole contour, whatever the tool can reach. If the tool cannot make a corner or a groove, the check reports a collision. Then turn the finishing pass off for this block, or use another tool.
\item \textbf{The check reports every collision.} It reports what the tool and its holder touch, also in rapid moves (Chapter~\ref{ch:checking}). Before the G-code is exported, WOP asks you if it found problems.
\end{itemize}

\begin{warningbox}
A block that has a problem it cannot plan (no contour, an error in the form) is \emph{left out} of the program, with its reason in the messages. Read the messages before you export. A program with a block left out is not what you meant.
\end{warningbox}

\section{Facing}\label{sec:blocks:facing}

\shot[0.45]{block_facing}{The form of a facing block.}

Facing cuts the end face of the part flat down to a Z that you give. The tool takes passes across the face. Each pass is one depth of cut deeper.

\begin{description}
\item[\field{Direction}] \emph{outside in}: the tool starts outside the part and moves to the centre. \emph{centre out}: it starts at the centre and moves out.
\item[\field{Target z}] The Z of the finished face.
\item[\field{Depth of cut}] How deep each pass is.
\item[\field{Overtravel past the centre (or bore)}] How far the tool goes past the centre so that no small point of material stays. For a tube, past the bore. The preset follows the tool (two nose radii).
\item[\field{Finish allowance} and \field{Finishing pass}] The last layer and its own pass (Section~\ref{sec:blocks:common}).
\end{description}

Each pass comes in beside the part, feeds the last millimetre into the material, cuts to the centre, and rises at feed until it is clear of the face. Only then does the tool move away in a rapid move.

\section{Turning}\label{sec:blocks:turning}

\shot[0.45]{block_turning}{The form of a turning block.}

The turning block makes a simple shape without a contour: a cylinder with a shoulder, a taper, a chamfer or a radius. Use it for the common jobs: roughing a bar to a diameter up to a Z. For a shape with several steps and arcs, draw a contour and use turning along a contour.

\begin{description}
\item[\field{Side}] \emph{outside} is turning. \emph{inside} is boring a bore; choose a boring bar as the tool.
\item[\field{Pass direction}] \emph{along Z}: the cuts run along the axis, one layer after the other. \emph{along X}: the cuts run across, from the outside toward the axis, one layer after the other.
\item[\field{Target radius} or \field{Target diameter}] The size that the surface will have. With X mode diameter, the label says diameter.
\item[\field{Start Z}] Where the cut begins. WOP takes the place where the material of this block begins: for example, the front of the part, or the step that an earlier block left. The value is shown grey and cannot be typed; the notes under the form say what was found.
\item[\field{End z}] Where the cut ends.
\item[\field{Taper}] The angle of a cone. \file{0} is straight. A positive angle makes the part larger towards the chuck. The target diameter is the one at Z~0. The button next to the field switches between the included angle in degrees and the rise of the diameter per 100~mm.
\item[\field{Start corner} and \field{Shoulder corner}] A chamfer or a fillet at the beginning of the cut and where the cut meets the shoulder. Choose the kind and give the size.
\end{description}

The fields \field{Depth of cut}, \field{Finish allowance}, \field{Finishing pass} and \field{Cutter compensation} are described in Section~\ref{sec:blocks:common}.

\section{Turning along a contour}\label{sec:blocks:contour}

\shot[0.45]{block_turning_contour}{The form of a block that turns along a contour.}

This block turns the part until a contour that you drew is left (Chapter~\ref{ch:contours}). It is the block for shapes that have several steps, arcs, grooves and tapers. It plans the roughing from the raw part, or from what the blocks before it left, down to the contour plus the finish allowance, and then the finishing pass along the contour.

\begin{description}
\item[\field{Contour}] The contour to cut. A block without a contour is left out of the program. A contour that cannot be used (an error in its drawing) is shown in the list with its reason.
\item[\field{Side}, \field{Pass direction}] As for turning (Section~\ref{sec:blocks:turning}). With a contour, \emph{along X} cuts layers across the part from the outside toward the axis, and after every cut the tool follows the contour a little way up and out.
\item[\field{Feed}, \field{Depth of cut}] The feed and the depth of cut of the roughing.
\item[\field{Finishing allowance}, \field{Finishing pass}, \field{Finishing feed}, \field{Spindle speed}] The finishing pass and its values.
\item[\field{Cutter compensation}] See Section~\ref{sec:blocks:common}.
\item[\field{Plunge feed} and \field{Plunge depth of cut}] Where the tool has to cut straight into the material, for example into a pocket behind a step, it can use another, usually smaller feed and depth. Empty: the values of the block.
\item[\field{Back clearance}] The smallest angle that the back of the tool keeps free from the contour (the free angle of the tool, in degrees). Where the back of the tool would touch the contour with less, roughing does not cut. The preset is 5~degrees.
\item[\field{Style}] The order in which the roughing is cut (Section~\ref{sec:blocks:styles}).
\end{description}

\begin{notebox}
Roughing stays off the finished part: it is cut only as far as the tool reaches, and the note under the block names what stays. But the finishing pass \emph{always} follows the contour, whatever the tool can reach. For a contour that your tool cannot make, for example a cavity behind a narrow mouth or an undercut, turn the finishing pass off. Otherwise the check reports a collision. Look at the picture and step through the block to see what the tool does.
\end{notebox}

\section{Drilling}\label{sec:blocks:drilling}

\shot[0.45]{block_drilling}{The form of a drilling block.}

Drilling makes a hole on the axis with a drill tool. The drill goes in along Z to an absolute depth.

\begin{description}
\item[\field{Drill mode}] \emph{drilling}: one stroke. \emph{peck drilling}: the drill goes in step by step and pulls back a little after each step to break the chips. \emph{peck drilling, full retract}: after each step the drill goes all the way out of the hole.
\item[\field{Depth (absolute Z)}] The Z that is reached. Behind the front face this number is negative.
\item[\field{Depth to full diameter}] Ticked: the depth is where the hole has its full diameter, and the point of the drill goes deeper. Not ticked: the tip of the drill reaches the depth.
\item[\field{Peck depth}] How deep each step goes (only for the peck modes).
\item[\field{Dwell time}] How long the drill waits at the bottom, in seconds.
\item[\field{Feed}] The feed of the drill in mm/rev.
\end{description}

The postprocessor can write a hole as a canned cycle of the controller (\file{G81}, \file{G82}, \file{G83}, \file{G73}) or as single moves. The LinuxCNC postprocessor that comes with WOP writes the cycles; the cycles were tried on the lathe, in the air and in material. If an earlier drill already cut part of the way, WOP skips the air.

\section{Threading}\label{sec:blocks:threading}

\shot[0.45]{block_threading}{The form of a threading block.}

The threading block cuts a thread along Z on the outside or inside of the part. The postprocessor writes the thread as one cycle of the controller (\file{G76} in LinuxCNC) if it has one. The fields are named like the words of that cycle.

\begin{description}
\item[\field{Side}] \emph{outside} or \emph{inside}.
\item[\field{Thread (M8x1, Tr40x7, 1/4-20 UNC)}] The designation of the thread. When you type it, WOP fills in the pitch and the sizes from it. You can also type the numbers yourself.
\item[\field{Spindle RPM}] The thread is cut at a fixed number of revolutions. A thread needs a constant spindle speed, not a constant surface speed.
\item[\field{Crest X radius}] The size in X (radius or diameter, as the label says) at the crest of the thread: for an outside thread its outer surface, for an inside thread the bore.
\item[\field{Pitch}] The distance between two turns of the thread.
\item[\field{Start Z (front)} and \field{End Z}] Where the thread starts and ends. The button \button{Abs.} / \button{Rel.} next to \field{End Z} lets you type the length of the thread instead of the Z of its end.
\item[\field{Thread depth}] Empty: WOP takes the depth of an ISO metric thread from the pitch.
\item[\field{First cut depth}, \field{Degression}, \field{Compound angle}, \field{Spring passes}] These set the passes. The notes under the fields explain them. The first pass is the deepest; with a \emph{Degression} of 2 every pass takes the same area, so the passes get shallower. A \emph{Compound angle} of about 29 or 30 degrees offsets each pass along the thread, so that one flank cuts more than the other. A \emph{Spring pass} repeats the last depth without a deeper cut.
\item[\field{Tapered at} and \field{Taper length}] A tapered run-in or run-out at the start or the end of the thread. \field{Taper length} is its length along the thread.
\end{description}

\begin{notebox}
WOP does not apply tolerance classes (6g, 6H). It cuts the thread to the nominal size. Trapezoid threads with several starts, UNJ threads and pipe threads (G, NPT) are not supported yet. Cut a thread in a test piece first and measure it before you cut a real one. The thread cycle of LinuxCNC was run on the lathe, also in material, and looked right; a finished thread has not been measured.
\end{notebox}

\section{Grooving}\label{sec:blocks:grooving}

\shot[0.45]{block_grooving}{The form of a grooving block.}

The grooving block cuts a groove with a grooving tool on the outside or the inside of the part, or parts the piece off.

\begin{description}
\item[\field{Style}] How the plunges are ordered (Section~\ref{sec:blocks:styles}).
\item[\field{Side}] \emph{outside} or \emph{inside}.
\item[\field{Groove Z (front edge)}] The Z of the wall of the groove that is nearer the front.
\item[\field{Width}] The width of the groove, along Z.
\item[\field{Final X radius} or diameter] The bottom of the groove. For a part-off, use \file{0}.
\item[\field{Peck depth} and \field{Retract after a peck}] A deep groove can be cut in steps. After each step the tool goes back a little (\emph{chip breaking}) or all the way out (\emph{full retract}).
\item[\field{One-sided cut in} and \field{Max one-sided cut}] A groove wider than the tool is cut in several plunges. A plunge may cut with one side of the blade free and the other side in material by at most this amount, in millimetres or as a percentage of the blade width. A smaller limit keeps the load on the blade lower.
\item[\field{Finish allowance} and \field{Finishing pass}] Finishing plunges at the walls and a pass along the bottom.
\end{description}

\begin{notebox}
A part-off needs a support for the piece that is cut off. WOP does not model that. When a block cuts the part into two pieces, WOP keeps the piece in the chuck and draws the other one dashed.
\end{notebox}

\section{Machine function}\label{sec:blocks:function}

\shot[0.45]{block_machine_function}{The form of a machine function block.}

A machine function block does not cut. It makes the machine do one thing: stop, wait, or stop the spindle. It has only a name and a function.

\begin{description}
\item[\emph{program stop}] The program stops until the operator presses start (\file{M0}). Use it to turn the part, to measure, or to clear chips.
\item[\emph{optional stop}] The same, but only if the optional stop switch of the control is on (\file{M1}).
\item[\emph{dwell}] The machine waits for a time (a field for the time in seconds appears).
\item[\emph{spindle stop}] The spindle stops.
\end{description}

The list can hold more entries: functions that your postprocessor defines for your controller. They have their own fields, which WOP shows when you choose them. Before the function, the tool goes to the safe position, so the operator can reach the part.

\section{Custom path}\label{sec:blocks:path}

\shot[0.45]{block_custom_path}{The form of a custom path block.}

The custom path is for moves that the other blocks cannot make. You draw the path of the tool yourself, and WOP checks it like any other block, draws it, and writes it. A path can have rapid moves, feed moves along lines and feed moves along arcs. Use it for a special groove, a thread relief, a chamfer on a strange shape, or a move to a place that the other blocks never go.

The form has the usual fields (name, tool, spindle, coolant), the \field{Feed} for pieces that have no feed of their own, and the button \button{Draw the path...}.

A new custom path block has nothing drawn yet. A block with no path is left out of the program, and the messages say so, until you draw a path.

\subsection{Drawing a path}
Click \button{Draw the path...}. The sketch editor opens (Chapter~\ref{ch:contours}). It shows two things that come from the program, in green:

\begin{itemize}
\item the \textbf{safe point}, a filled green square with its name. This is where every block starts and where the tool waits. WOP takes it from the Setup (Section~\ref{sec:project:safe});
\item the \textbf{part as it is before this block}: the material that the blocks before it left. It is drawn as green lines and arcs. You cannot move them.
\end{itemize}

\shot{sketch_path}{The sketch editor for a custom path. The green square is the safe point, the green lines are the part before the block, the pink line with the big arrow is the piece that is shown now while stepping.}

Draw the path as a chain of lines and arcs with the usual tools (Section~\ref{sec:sketch:draw}).

\begin{itemize}
\item \textbf{The first line starts at the safe point.} Click the green square as the start. If the path does not start there, the check under the drawing says so.
\item \textbf{The pieces follow each other in the order you draw them}, each one starting where the one before ended. This order is the order of the moves.
\item \textbf{Every line is a feed move unless you mark it as rapid.} Select the line, or click the button, and use \emph{Rapid move on / off} in the section \emph{Mark}. A rapid line is drawn dashed. A rapid move has no feed: the control decides how fast it is. An arc is always a feed move.
\item \textbf{A piece can have a feed of its own.} Select it and type the feed under the drawing (mm/rev). Empty means the feed of the block. A rapid line has no feed field.
\item \textbf{A path may go back along its own line and may cross itself.} Draw the second line over the first one. They are separate pieces and keep their order. Often you cut in with a feed line and go back out with a rapid line along the same line.
\item A fillet or a chamfer between two lines keeps the flags of the lines: a corner between two rapid lines is rapid. If both lines have the same feed of their own, the corner gets that feed; if the feeds differ, the lines keep theirs and the corner uses the feed of the block.
\end{itemize}

You can end the path anywhere you want. After the block, WOP moves the tool to the safe point (or, before another tool, to the tool change position), and the check reports it if that move is blocked by material.

\begin{notebox}
While the tool feeds, the material it passes through counts as removed. A path may start inside the material. Check what you drew with \button{Step through} in the main window (Section~\ref{sec:window:stepping}).
\end{notebox}

\subsection{Dimensions to the part}
You can dimension the path to the green part, so that the path follows when the part changes. Choose the \emph{Dimension} tool, click a point of your path, then click a green line, and place the number: you get the \emph{distance} from the point to the line, never a diameter. You can also dimension to the points of the part.

\begin{warningbox}
Use as few dimensions to the part as you can. They are easier to redo when the part changes. The green part is the material that the blocks before this block leave. When an earlier block changes it, WOP looks at the outline of the part. If it still has the same pieces (the same number and kinds of lines and arcs) and every dimension still works, the dimensions stay and the path follows. If the number or the kind of the pieces has changed, \emph{every} dimension that touches the part is dropped, and the dialog tells you how many. A dimension that cannot be met any more is dropped too. The path itself stays where it is.
\end{warningbox}

\subsection{Is the path fixed?}
A path that you have not fully dimensioned can still move its points. The block adds a note \emph{The path is not fully constrained} to the messages whenever this is so. It is only a note: the path is planned and cut as it is. Add dimensions and relations when you want the path to follow changes in the part. You do not have to.

\subsection{Stepping through the path}\label{sec:blocks:pathstep}
A path can come back along its own lines, so it can be hard to see in what order the moves are made. The sketch editor can step through the path piece by piece. Use the buttons \button{< Step back} and \button{Step forward >} under the drawing, or the keys \key{Left} and \key{Right}.

\begin{itemize}
\item The first piece is shown and selected, with a big arrow at its end that shows the way the tool moves. All later pieces are hidden.
\item Each step shows one more piece. The older pieces are drawn thin and faint, and only the dimensions of the piece you are on are shown.
\item The piece you are on is the one you select and change with the mouse, also where several pieces lie on one place. Hidden pieces cannot be picked.
\item A label says which piece it is and what kind: \emph{Step 3 of 7: feed line}.
\item If a piece does not follow the chain, stepping stops there, and the message says why.
\item \key{Esc} (when no tool is on) or stepping back from the first piece shows everything again. Starting a drawing tool also ends the stepping, because a new piece is added at the end of the path.
\end{itemize}

\section{Styles}\label{sec:blocks:styles}

Some blocks have a \field{Style}. A style decides only the \emph{order} and the \emph{depth} of the steps. It never changes what is cut. The style is part of the program; there are no style files to edit.

Open the list of styles and hold the mouse over an entry: a window shows pictures of the style on small samples. The colour in the pictures goes from blue (made first) to red (made last), and only the cutting moves are shown.

\subsection{Grooving}
\begin{description}
\item[Standard] Plunges from the front wall towards the chuck, each plunge at most one one-sided cut wider than the one before. Finishing plunges at the walls and a pass along the bottom.
\item[Layers] Cuts the groove layer by layer: first all plunges to a depth, then the next layer deeper.
\end{description}

\shot[0.6]{style_grooving_standard}{Grooving style \emph{Standard} on four samples.}
\shot[0.6]{style_grooving_layers}{Grooving style \emph{Layers}.}

\subsection{Turning along a contour}
\begin{description}
\item[Standard] Every roughing line is cut when it is found, level by level from the front. Pockets and open lines are taken together.
\item[Pockets last] All open lines are cut first. Then the pockets behind crests, the front one first, layer by layer.
\end{description}

\shot[0.6]{style_contour_standard}{Contour style \emph{Standard}.}
\shot[0.6]{style_contour_pockets_last}{Contour style \emph{Pockets last}.}

\chapter{Contours and the sketch editor}\label{ch:contours}

Many things in WOP are shapes: the finished outline of a part, the outline of a tool, the outline of jaws, the shape of a raw part, and the path of a tool. You draw all of them in the same place, the \emph{sketch editor}. This chapter explains how to work in it. It starts with contours, which are the shapes of parts, because they are the most common use.

\section{Contours in the flow list}\label{sec:sketch:contours}

A \emph{contour} is the outline of a part, drawn as a half section: X is the radius and goes up, Z is the length and goes to the right, and the spindle axis is a dash-dotted line. A block that works along a contour (turning along a contour, Section~\ref{sec:blocks:contour}) cuts the material away until the contour is left.

A contour lives under the heading \menu{Contours} in the flow list. It is not a step of the program. Many blocks can use the same contour: a roughing block and a finishing block can both use the shape of the finished part.

\subsection{Making a contour}
Choose \button{Add block}, then \menu{Contour}. A new contour is made, and the sketch editor opens for it at once. Draw the outline and press \button{OK}. The contour is now in the list.

\subsection{The form of a contour}
Click a contour in the flow list. Its form has the name of the contour and the button \button{Sketch editor...}. The text under the name says \emph{The contour is fine}, or what is wrong with it. Give the contour a name that says what it is, for example \emph{stepped shaft}. A block chooses its contour by this name.

\shot[0.5]{contour_form}{The form of a contour: its name and the button that opens the sketch editor.}

\subsection{Rules for the contour of a part}
The editor tells you under the drawing what is wrong. A contour of a part must be

\begin{itemize}
\item \textbf{one open chain} of lines and arcs that follow one another. If the outline is in several parts, join them, or turn the extra ones into construction lines (Section~\ref{sec:sketch:marks});
\item \textbf{open}, not a closed loop. Close the shape with the spindle axis, not with a line;
\item \textbf{above the spindle axis}: no point may have an X below zero;
\item \textbf{free of crossings}: no piece may cross or touch another piece, except where two pieces join.
\end{itemize}

You may draw the contour from the front to the chuck or the other way round. WOP turns it so that it starts at its highest Z.

\subsection{Deleting and copying}
\button{Delete} removes the chosen contour. If blocks use it, WOP asks whether to delete these blocks too, or to keep them without a contour. A block without a contour is left out of the program until you give it one. A right click on a contour and \menu{Copy} makes a copy that you can change without changing the original.

\section{The sketch editor window}\label{sec:sketch:window}

The sketch editor opens in its own window. It opens

\begin{itemize}
\item from the button \button{Sketch editor...} in the form of a contour;
\item by itself when you add a new contour;
\item from the buttons in the tool view and the workholding view, to draw the shape of a tool or of jaws (Chapters~\ref{ch:tools} and~\ref{ch:workholding});
\item from \button{Draw the raw part...} in the Setup (Section~\ref{sec:project:draw});
\item from \button{Draw the path...} in a custom path block (Section~\ref{sec:blocks:path}).
\end{itemize}

There is no menu entry for it. The editor knows what it was opened for, and it checks the drawing against the rules of that use. The title of the window says what it is.

\shot{sketch_editor}{The sketch editor with a contour. On the left the buttons in five sections, in the middle the drawing, under it the hint and the check of the outline.}

\subsection{The parts of the window}
\begin{description}
\item[The buttons] on the left are square with a small picture. Hold the mouse over a button to see what it does (Figure~\ref{fig:sketch_hover_text}). They are in five sections: \emph{Draw}, \emph{Edit}, \emph{Constraints}, \emph{Mark} and \emph{General}.
\item[The drawing] fills the middle. It shows a grid, the origin (a small ring) and the spindle axis. Use the mouse wheel to zoom, and drag with the middle mouse button to move the picture.
\item[The hint] under the drawing says what to click next for the tool that is on.
\item[The check] under the hint says \emph{The outline can be used}, or lists what is wrong. \button{OK} stays off until the list is empty.
\item[The tray] on the right of the drawing lists the relations and dimensions (Section~\ref{sec:sketch:tray}). It starts folded; click its handle to open it.
\end{description}

\shot[0.7]{sketch_hover_text}{The hover text of a button.}

\subsection{OK, Cancel and Undo}
\button{OK} puts the drawing into the program. \button{Cancel} throws the changes of this visit away. The key \key{Esc} never closes the window: it ends what you are doing (a line you are drawing, a tool that is on) and goes back to selecting. Undo and redo of the editor work with \key{Ctrl}+\key{Z} and \key{Ctrl}+\key{Y}, and with the buttons in the section \emph{General}. They only concern this window; the main window has its own undo (Chapter~\ref{ch:edit}).

\section{Drawing}\label{sec:sketch:draw}

Choose a tool in the section \emph{Draw}. The tool stays on until you click its button again, press \key{Esc}, or choose another tool. Always read the hint under the drawing: it tells you the next click.

\subsection{Lines}
Click the start, then the end. The next line continues from the end of the last one, so you draw an outline in one go. A right click ends the line; \key{Esc} ends it and goes back to selecting. With \button{Snap} on (the default), a line that is nearly horizontal or vertical is made exactly so, and the points snap to existing points.

\subsection{Circles and arcs}
\begin{description}
\item[Circle] Click the centre, then a point on the circle.
\item[Arc (centre, start, end)] Click the centre, then the start, then the end of the arc. The end sets the direction.
\item[Arc (start, end, point on arc)] Click the start, then the end, then a point the arc goes through.
\end{description}

\subsection{Points}
The \emph{Point} tool puts a point where you click. A point on its own does not belong to the outline. You need it for a name (Section~\ref{sec:sketch:names}) or as a place to dimension to.

\subsection{Snapping}
The button \button{Snap} in the section \emph{General} turns snapping on and off. With snapping, a click near a point uses that point, so two lines that meet really share a point. This matters: lines that only look joined but do not share a point make an outline in two parts, and the check says so.

\section{Selecting and moving}\label{sec:sketch:select}

When no tool is on, you are selecting. Click a point or a piece to select it; it turns pink. Hold \key{Ctrl} to add to the selection or to take away from it. Drag on empty space to select with a box. \key{Delete} removes the selection and every piece that ends on a selected point.

Drag a selected point or piece to move it. What happens depends on what is already fixed. A free point follows the mouse. A point that is tied by relations or dimensions moves only as far as they allow; if the mouse asks for more, the point follows a part of the way. If you drop a point on another point, the two become one.

\section{Relations}\label{sec:sketch:relations}

A \emph{relation} is a rule that you give to the drawing, for example \emph{this line is horizontal}. WOP keeps the rule when you move something else. Relations are what make a drawing stay the way you drew it. They are in the section \emph{Constraints}.

Click the button of a relation. If the right things are selected, it acts on them at once. If not, the button stays pressed and you click the things it needs, as the hint says. Click the button again or press \key{Esc} to stop.

\begin{center}
\begin{tabularx}{\textwidth}{@{}l l Y@{}}
\toprule
\textbf{Relation} & \textbf{Tag} & \textbf{What it does} \\
\midrule
\menu{Horizontal} & H & Two points, or the ends of a line, get the same X: the line runs along Z. The origin can be one of the two points. \\
\menu{Vertical} & V & Two points, or the ends of a line, get the same Z: the line runs across Z. \\
\menu{Parallel} & $\parallel$ & Two lines run in the same direction. \\
\menu{Perpendicular} & T & Two lines meet at a right angle. \\
\menu{Equal} & = & Two lines get the same length, or two arcs or circles the same radius. \\
\menu{Tangent} & tan & A line and an arc, or two arcs, that meet at a point touch without a bend. \\
\menu{Fix in place} & F & A point or a piece cannot move any more. \\
\menu{Point on line / arc} & on & A point lies on a line or an arc. \\
\menu{Midpoint} & mid & A point is in the middle of a line. \\
\menu{Join points (coincident)} & & Two points become one. \\
\bottomrule
\end{tabularx}
\end{center}

The tags are small letters that the drawing shows next to the relation. You can hide them with the button \emph{Show the relations}. A relation that cannot hold is refused with a message, and nothing changes.

\section{Dimensions}\label{sec:sketch:dimensions}

A \emph{dimension} is a number that you give to the drawing: a length, a distance, a radius, an angle. The shape follows the number. A dimension is also a relation, but it has a value.

\subsection{Placing a dimension}\label{sec:sketch:placing}
Choose the \emph{Dimension} tool, the first button of the section \emph{Constraints}. Then

\begin{itemize}
\item click \textbf{a line} and then where the number should stand. Near the line it is the \emph{length}. To the side of the line it is the distance along X; above or below it, along Z;
\item click \textbf{two points} (the origin too) and then where the number should stand: the \emph{distance} between them, in the same way;
\item click \textbf{an arc} for its radius, or \textbf{a circle} for its diameter;
\item click \textbf{two lines} for the \emph{angle} between them;
\item click \textbf{a point or a line and a centre line} for a radius, or for a diameter (Section~\ref{sec:sketch:centre}).
\end{itemize}

The number is shown while you move the mouse. A click places it. The dimension is made, and the drawing moves to it.

\subsection{Changing and moving a dimension}
Double click a number, or choose the dimension in the tray and edit it there, to type another value. A point and a comma both work as the decimal point. A dimension that cannot be met together with the others is refused with a message, and the old value stays. You can drag the lines of a dimension with the mouse to move the number to a better place; this changes nothing in the shape.

\shot{sketch_editor_tray}{The tray lists the relations and dimensions. A chosen entry is drawn in pink in the picture.}

\subsection{The tray}\label{sec:sketch:tray}
The tray is under the heading \emph{Relations and dimensions}. Each relation and each dimension is a line in it. Click a line to see what it acts on in the picture, and \button{Remove} to delete it. The line at the bottom of the tray says \emph{Fully defined}, or how many free movements are left (Section~\ref{sec:sketch:defined}).

\subsection{Diameters in a project typed in diameters}
If the X mode of the project is \emph{diameter} (Section~\ref{sec:project:xmode}), a dimension along X is shown and typed as a diameter, because a lathe worker thinks that way. The drawing itself stays in radii.

\section{Fully defined or not}\label{sec:sketch:defined}

In a sketch a piece is \textcolor{blue!60!black}{blue} while it can still move, and black (white with a dark theme) when it cannot move any more because relations and dimensions hold it. The text in the tray counts the movements that are left. You do not have to make every sketch fully defined. But you can: when the tray says \emph{Fully defined}, the sketch is exact and nothing can move by a slip of the mouse.

\begin{tipbox}
Dimension and constrain a sketch \emph{inside itself}: put relations and dimensions between its own parts (this line is as long as that, this point is as far from that point). Fix its position to the origin \emph{last}. Then you can change the position or the direction of the whole shape later with one change, and the shape itself stays as you made it.
\end{tipbox}

\section{Centre lines and diameters}\label{sec:sketch:centre}

A \emph{centre line} is a helper line that dimensions are measured \emph{across}. In a drawing for a turning contour the spindle axis is a centre line that is always there; you do not draw it. Click a point or a line and then the spindle axis with the \emph{Dimension} tool: if you place the number between the thing and the axis (or beyond the thing), you get a \emph{radius}, the distance to the axis. If you place it across the axis, you get a \emph{diameter}, the distance to the mirror image of the point on the other side. The mirror image is made for you.

You can turn any other line into a centre line with the button \emph{Centre line on / off} in the section \emph{Mark}. Use this for a shape that has its own axis, for example a hole.

\section{Trim, mirror, fillet and chamfer}\label{sec:sketch:edit}

The section \emph{Edit} changes pieces that are already there.

\begin{description}
\item[Trim] Click the part of a line or arc that you want to cut away. It is cut up to the next crossing or to the end. Hold the button and drag across several pieces to trim everything the mouse crosses.
\item[Mirror] First select the pieces you want to mirror. Then click the \emph{Mirror} button and click the line to mirror about: any line, a construction line or centre line too, or the spindle axis. The mirrored copies are made and selected.
\item[Fillet corner...] Rounds a corner with an arc. WOP asks for the radius when you press the button; every corner you click after that gets this radius. A corner is a point where exactly two lines or arcs meet. Fillets also work between a line and an arc and between two arcs.
\item[Chamfer corner...] The same, but cuts the corner with a straight line, a given distance from the corner along each piece.
\end{description}

A fillet or a chamfer is one unit. The two pieces are shortened, the new arc or line is added, and the sharp corner stays as thin dashed lines (\emph{reference lines}). If you change the radius, or delete one part of the unit, the unit is undone and the sharp corner comes back.

\section{Marks: construction lines, cutting edges and names}\label{sec:sketch:marks}

The section \emph{Mark} says what a line or a point \emph{is}, without changing where it is.

\begin{description}
\item[Construction on / off] A \emph{construction line} (or arc) is a helper. It is drawn dashed, and it is not part of the outline. Use it to hold a dimension or to place a point. Select pieces and click the button to make them construction lines, or real again.
\item[Centre line on / off] See Section~\ref{sec:sketch:centre}.
\item[Cutting edge on / off] Only in the sketch of a tool: marks the pieces of the outline that cut (Chapter~\ref{ch:tools}).
\item[Rapid move on / off] Only in the sketch of a custom path (Section~\ref{sec:blocks:path}).
\end{description}

\subsection{Named points}\label{sec:sketch:names}
Some shapes need points with a special meaning. A named point has a ring and its name next to it. In the section \emph{Mark} there is one button for each kind of name that the shape you are drawing can use. The table shows the names. Select a point and click the button, or click the button and then the point. Click again to take the name away.

\begin{center}
\begin{tabularx}{\textwidth}{@{}l l Y@{}}
\toprule
\textbf{Shape} & \textbf{Name} & \textbf{Meaning} \\
\midrule
Raw part & Clamping plane & Where the jaws begin (only its Z). Section~\ref{sec:project:draw}. \\
Tool & Tip & The touch-off point of the tool. It has no button: WOP finds it from the shape unless you name it. Chapter~\ref{ch:tools}. \\
Jaws & Mount, Step 1, Step 2 \ldots & Where the jaws meet the chuck face, and the steps where a part is gripped. Chapter~\ref{ch:workholding}. \\
Chuck body & Face & The front face of the chuck. Chapter~\ref{ch:workholding}. \\
\bottomrule
\end{tabularx}
\end{center}

\section{Importing a DXF drawing}\label{sec:sketch:dxf}

If you have the shape as a drawing from a CAD program, the button \emph{Import DXF} in the section \emph{General} loads it. It replaces what is in the sketch, and WOP asks first if there is something in it.

For a contour of a part WOP asks \emph{How does the part lie in the drawing?}: with the spindle axis horizontal (drawing X is Z, drawing Y is the radius), or vertical (drawing Y is Z, drawing X is the radius). For tool and workholding shapes, the drawing is taken as it is, with the spindle axis horizontal; nothing is moved, turned or mirrored, and the origin of the drawing stays the origin. A \emph{point} with a \emph{text} next to it in the drawing becomes a named point.

WOP reads lines, arcs and polylines (a polyline is split into its lines and arcs). It does not read circles, ellipses, splines or dimensions of the drawing. After the import, a message tells you what was taken and what was left out. The imported shape is a normal sketch: you can dimension it and change it.

\section{Locked things}\label{sec:sketch:locked}

In some sketches parts of the drawing are \textcolor{green!45!black}{green}. They come from the program, not from you: the safe point and the part before a block in the sketch of a custom path. You can pick them and dimension to them, but you cannot move, change or delete them. WOP tells you so if you try. Section~\ref{sec:blocks:path} explains how to use them.

\section{What the editor does not do}\label{sec:sketch:limits}

\begin{itemize}
\item It cannot offset a contour (make a parallel copy at a distance) or turn a whole shape by an angle. Draw the line you need, or select everything and move it.
\item It does not read circles, ellipses and splines from a DXF file.
\item A dimension that only shows a size and does not drive the shape (a \emph{reference dimension}) does not exist yet.
\end{itemize}

\chapter{Tools}\label{ch:tools}

In WOP a tool is not a number in a table. It is a shape with a cutting edge. WOP uses the shape to plan the cuts and to find collisions of the tool, its holder or the bar with the part. This chapter says how to make a tool, how its shape is made, and how the tool library works. The menu is in Chapter~\ref{ch:tools_menu}.

\section{The Tools view}\label{sec:tools:view}
Open it with \menupath{Tools}{Show tools}. The top row has the toolbox of the project (see below) and the buttons. Under it are the list of the tools, a picture of the chosen one, and a few lines of information about it. The form of the chosen tool is on the right; a change in the form is used at once.

\begin{description}
\item[\button{Add tool}] Adds a tool of a type (next section).
\item[\button{Add from library}] Copies a tool of the tool library into the project.
\item[\button{Save to library}] Writes the chosen tool to the tool library in use. If the library has another tool of the same name, WOP asks whether to replace it. The button then reads \button{Update library tool}.
\item[\button{Save as toolbox...}] Makes a toolbox of the tools of the project (Section~\ref{sec:tools:toolboxes}).
\item[\button{Delete}] Removes the chosen tool from the project.
\item[\button{Sketch editor...}] Opens the shape of the tool in the sketch editor.
\item[\field{Tool offset}] The offset number that the postprocessor calls after the tool change (for LinuxCNC \file{G43 H...}). With \emph{tool number}, the offset has the number of the tool. Set another number if the offset of this tool sits elsewhere in your tool table.
\end{description}

In the list, a tool reads like \emph{T3  Grooving tool  (groove, steel)}. A star after the name means that the tool differs from the tool of the same name in the library.

\section{Tool types}\label{sec:tools:types}
Every type has a form with the numbers of its shape. The form always starts with the \field{Number}, the \field{Name} and the \field{Material} of the cutting edge. The hand (right or left) mirrors the tool.

\begin{description}
\item[Turning tool] An insert tool for the outside. The fields are the \field{Nose radius}, the lengths and angles of the side edge and the front edge, and the width and length of the tool. The angle between the edges at the tip is shown as information.
\item[Turning tool (ISO holder)] Made from the codes of a standard holder (for example \emph{PCLNR 2525 M12}) and a standard insert (for example \emph{CNMG 120408}). Choose or type the \field{Holder code} and the \field{Insert code}; the shape, style, section, length, insert size and corner radius are worked out and can be changed.
\item[Boring bar] A tool for the inside. Like the turning tool, with the \field{Wall edge}, the \field{End edge}, the front width, the width of the bar and its length. The information shows the smallest bore this bar fits in.
\item[Boring bar (ISO)] Made from the \field{Boring bar code} and the insert code, with the diameter and the length of the bar.
\item[Grooving tool] A blade with a width, a reach, corner radii and a holder. It can work outside or, with \field{Inside}, in a bore. \field{Touch-off corner} says which corner of the blade is the zero point of the tool.
\item[Threading tool] A tip with an \field{Included angle}, a lead angle, a tip radius and a flat, and the size of its body. It can work outside or inside.
\item[Drill] The \field{Diameter}, the \field{Length} and the \field{Tip angle}.
\end{description}

\begin{notebox}
For the ISO holders and bars, WOP draws the head of the holder as a plain envelope; the real head is not standardised. A few dimensions are guessed from the standard and are marked as such in the information under the picture. Compare the picture with your real tool before you trust the collision check for it.
\end{notebox}

\section{The shape of a tool}\label{sec:tools:shape}
The shape comes from the numbers of the form, so changing a number redraws the tool. If the form is not enough, \button{Sketch editor...} opens the shape for drawing. After you change it there, the tool becomes a \emph{free shape}: it is no longer made from the numbers of its type, and the form shows only the name, number and material.

In the sketch editor (Chapter~\ref{ch:contours}) the shape must be closed. The tool is drawn so that the tip, the point where you touch off, is the zero point. Two marks matter:
\begin{itemize}
  \item \textbf{Cutting edges} are drawn in orange. Only a cutting edge may touch material. Everything else of the tool must never touch it. A line becomes a cutting edge with the mark \emph{cutting edge} (Section~\ref{sec:sketch:marks}).
  \item \textbf{Named points} are shown with their name.
\end{itemize}

\section{Named points}\label{sec:tools:points}
A tool shape can carry names for points, such as the tip. They are drawn with their name in the sketch editor and are used for the touch-off point and the checks. The tools made from a form have their points already; you only need to think about them when you draw a tool by hand.

\section{The tool library}\label{sec:tools:library}
The tool library keeps tools you use again and again. Make a tool once in a project, press \button{Save to library}, and use \button{Add from library} in the next project.

A project holds its own copy of its tools. WOP remembers which library tool a project tool came from. When you close the project, WOP compares: if a tool of the project differs from the library, or is not in it yet, it asks whether to update the library (Section~\ref{sec:file:questions}). The library you use is chosen under \menupath{Tools}{Library}.

\section{Toolboxes}\label{sec:tools:toolboxes}
A toolbox is a list of library tools with the numbers they have in the turret of your lathe. A new project starts with the tools of the \emph{default toolbox}. Without a default toolbox it starts with no tools.

\begin{description}
\item[In the Tools view] \button{Save as toolbox...} makes a toolbox of the tools of the project. The tools must be in the library first. If the tools came from a toolbox, you can update it or make a new one. \button{Make default} makes the toolbox of the project the default.
\item[\menupath{Tools}{Edit toolboxes...}] The window has the list of the toolboxes on the left with \button{New}, \button{Rename}, \button{Delete} and \button{Make default}, and the tools of the chosen toolbox on the right with \button{Add tool} and \button{Remove tool}. You can set the number of each tool. The tools you add come from the library.
\item[\menupath{Tools}{Use toolbox}] Puts a toolbox into the project (Chapter~\ref{ch:tools_menu}).
\end{description}

\chapter{Workholding}\label{ch:workholding}

Workholding is what holds the part: the \emph{chuck body} and the \emph{jaws}. WOP draws them behind the part and uses their shapes for the checks: a cut that runs into the jaws is reported. The menu is in Chapter~\ref{ch:workholding_menu}; the choice for a project is made in the Setup (Section~\ref{sec:project:holding}).

\section{Chuck and jaws}
There are two types, both in the workholding library:
\begin{description}
\item[Chuck body] The fixed body of the chuck. It has a \field{Name}.
\item[Chuck jaws] The jaws that grip the part. The fields are the \field{Name}, the \field{Number of jaws} and the \field{Width of one jaw} in degrees (they are used for the 3D picture), and the \field{Chuck body} that belongs to them. If you choose jaws in the Setup, their chuck body is used too.
\end{description}

\section{The shape}\label{sec:workholding:shape}
A new item has a placeholder shape. Draw the real one with \button{Sketch editor...} (Chapter~\ref{ch:contours}). Draw it as it sits on the machine:
\begin{itemize}
  \item X is the radius and Z runs along the axis.
  \item The front face is at $Z=0$ and the body goes towards minus Z.
  \item The shape must be closed.
\end{itemize}
The named points tell WOP how the parts fit:
\begin{itemize}
  \item On jaws: \emph{Step} points mark the steps that can sit on the part, and \emph{mount} marks where the chuck body meets the jaws.
  \item On a chuck body: \emph{face}.
\end{itemize}
Name the points with the buttons of the section \emph{Mark} in the sketch editor (Section~\ref{sec:sketch:names}).

\section{Placing the part}
In the Setup you choose the jaws and, for stepped jaws, the field \field{Jaw point on the part}: the step of the jaws that sits on the part. The \field{Stick-out} says how far the part comes out of the jaws. With \field{Grip inside} the jaws grip a tube from the bore (Section~\ref{sec:project:holding}).

With no workholding chosen, nothing is drawn and the checks use a flat clamping plane at the stick-out.

\section{What is checked}
\begin{itemize}
  \item A cut that touches the jaws is reported as a collision (\emph{touches the jaws}).
  \item A cut that goes behind the clamping plane of the Setup is reported too.
  \item The Setup also warns when the stick-out is longer than the part, when a bar is too thick for the chuck body, or when the jaws of an inside grip reach into the wall of the tube.
\end{itemize}
WOP does not check how far a part may stick out before it bends or vibrates.

\begin{notebox}
Stepped jaws and the other workholding checks are only as good as the shape you draw. Compare the picture with your real chuck.
\end{notebox}

\chapter{Machines}\label{ch:machines}

A machine in WOP holds what belongs to your lathe and not to one part: the travel of the axes, the speeds, the power, the tool change position and the postprocessor. Make a machine once; every project can then use it.

Open the list with \menupath{File}{Machines...} (Figure~\ref{fig:machines}).

\shot[0.85]{machines}{The window \emph{Machines}.}

The machines are on the left. The one that the project uses is marked \emph{(in project)}, the default machine \emph{(default)}. The form of the chosen machine is on the right. A change is saved when you leave the field. A value that makes no sense stays in the field with a message below the form until you correct it.

\section{The buttons}
\begin{description}
\item[\button{New}] Adds a machine with preset values.
\item[\button{Import...}] Makes a machine from a file. WOP asks what kind of file it is: a \emph{LinuxCNC ini file} or a \emph{WOP machine file}. The values that the file gave are marked as set; the others stay preset. A message lists what was assumed and what you still have to enter.
\item[\button{Export...}] Writes the chosen machine to a WOP machine file (\file{.machine.json}).
\item[\button{Delete}] Removes the machine from the list. Projects keep their own copy.
\item[\button{Make default}] The default machine is used for new projects, and for old projects that have no machine.
\item[\button{Use in project}] The project uses the chosen machine.
\item[\button{Add project's machine}] Adds the machine of the open project to your list, when it is not there yet.
\end{description}

\section{The settings of a machine}\label{sec:machines:settings}
\begin{description}
\item[\field{Name}] The name in the lists.
\item[\field{Unit system} and \field{Default x mode}] Millimetres, and whether X is shown as diameter or radius in new projects of this machine.
\item[\field{X min}, \field{X max}, \field{Z min}, \field{Z max}] The travel of the axes in mm. WOP compares the size of the part and of the cuts with them. X is a radius.
\item[\field{Spindle min} and \field{Spindle max}] The speed range in RPM.
\item[\field{Spindle power}] In kW. The power check compares the power a cut needs with this value.
\item[\field{X max velocity} and \field{Z max velocity}] The fastest feed move of each axis, in mm/min. If a feed per revolution at the chosen speed is faster than an axis can go, the messages say so.
\item[\field{Constant surface speed}] Tick it if the controller can keep the cutting speed constant (G96 on LinuxCNC). If it is off and a block asks for constant surface speed, the messages say so.
\item[\field{Spindle synchronised moves}] Tick it if the controller can move in step with the spindle. Threading needs it.
\item[\field{Postprocessor}] The postprocessor that writes the G-code for this machine (Chapter~\ref{ch:posts}).
\item[\field{Tool change position, X radius} and \field{Z}] Where the tool is changed, in the coordinates of the part. If you leave them empty, the position of the Setup is used.
\item[\field{Tool change in machine coordinates, X radius} and \field{Z}] A fixed place of the tool change in machine coordinates. If you set it, the postprocessor first moves there with a machine-coordinate move (\file{G53} for LinuxCNC) before every tool change, whatever the work offset is. If you leave it empty, no such move is written.
\end{description}

A field that you did not touch has a lighter colour. A right click on a field opens \menu{Reset to default}.

\section{Making a machine for your lathe}\label{sec:machines:own}
\begin{enumerate}
  \item Open \menupath{File}{Machines...} and press \button{Import...} if you have the \file{.ini} file of your LinuxCNC configuration. The travel, speeds and similar values come from it.
  \item Otherwise press \button{New} and fill in the form. The five values WOP needs are the limits of X and Z, the minimum and maximum RPM and the spindle power. The import lists the ones that are still missing. A power or RPM check can only compare with a value that is set.
  \item Choose the postprocessor, and the tool change position if you have a fixed one.
  \item Press \button{Make default} if this is the machine you use most, and \button{Use in project} for the open project.
\end{enumerate}

\begin{warningbox}
WOP uses the limits to warn you. It does not stop a program that is outside them: the G-code is written as drawn. The real limits and the safe position are still your job on the machine.
\end{warningbox}

\chapter{Postprocessors}\label{ch:posts}

\begin{warningbox}
Every postprocessor can crash a machine, the ones that come with WOP too. None of them was run on every controller. Read the G-code, and cut in air first, before you run a new postprocessor on a part.
\end{warningbox}

\section{What a postprocessor does}
WOP plans the paths in its own form: moves, spindle, coolant, tool changes. A \emph{postprocessor} (post) turns them into the G-code of your controller. It decides how a rapid move is written, whether X is a diameter, how a thread or a drill cycle looks, what the first lines say, and so on.

A post is a text file that ends in \file{.post.ini}. You do not need a programmer to change it. You change it in the postprocessor editor, which shows what the G-code will look like while you edit.

One post comes with WOP: the post for \textbf{LinuxCNC}. It writes X as diameter, feed per revolution, constant surface speed, the thread cycle and the drill cycles of LinuxCNC, and cutter compensation (G41, G42, G40) for a finishing pass that asks for it. It was tried on one lathe.

Inside the program there is also a plain default post (plain G0, G1, G2, G3 and a few M codes, feed in mm/min, no cycles). It is not in the lists. A post that does not have a page of its own for something uses the page of this default post.

The folder \file{examples/posts} holds reference posts to import with \button{Import...}: a Fanuc style post, a plain G-code post, and the LinuxCNC post with a custom machine function (a bar feeder) that shows how a custom function is written. The Fanuc and the plain post are written from a rough understanding of these controllers and are \emph{not tried} on any machine. Use them as a start for your own post.

The machine says which post is used (Chapter~\ref{ch:machines}).

\section{The editor}\label{sec:posts:editor}
Open it with \menupath{File}{Postprocessors...} (Figure~\ref{fig:post_editor}).

\shot[0.9]{post_editor}{The postprocessor editor.}

On top is the list of the posts and the buttons:
\begin{description}
\item[\button{New}] Makes a new post.
\item[\button{Save}] Saves your changes.
\item[\button{Save as copy...}] Saves under a new name. The posts that come with WOP cannot be changed; make a copy and change that.
\item[\button{Import...}] Copies a post file of somebody else into your posts.
\item[\button{Delete}] Removes one of your posts.
\item[\button{Folder}] Shows the folder of your posts in Windows.
\item[\field{Sample X}] Whether the sample on the pages shows X as diameter or radius.
\end{description}
The title of the window gets a star when there are unsaved changes.

\section{The pages of a post}\label{sec:posts:pages}
The pages are on the left. They are in this order.

\begin{description}
\item[Post] The name and the description of the post, and the name of its file. A note says whether the whole file reads without a problem.
\item[Settings] Numbers and switches: the number of decimals, line numbers, whether X is a diameter, whether arcs can be written while X is a diameter, how a thread and a drill are written (as cycle, in two blocks, or as plain passes), cutter compensation, a speed ramp before the spindle stops (for a drive that trips when it is stopped from speed), the characters a comment may not hold, and the file extension. Hover over a name to see what it means.
\item[Code table] The G and M codes of the controller, by name. WOP uses the names; you set the codes. A code that is empty means the controller has no such function, so WOP does not write it. The section \emph{Work offsets} is the list of work offsets of the controller (G54, G55 \ldots) that the Setup offers; add the ones your controller has. You can add codes of your own.
\item[Modal groups] Which codes cancel each other (for example G0, G1, G2 and G3 are one group). WOP uses the groups to leave out a code that is still valid from the line before. List the codes of the code table here; a code of the table that is in no group is named on the Code table page. You can add groups of your own.
\item[The concepts] After these four pages comes one page for each thing the post writes: the header and the footer, comments, the start of a block, rapid, the move to the tool change position (a page of its own, for a controller that has a code for it, such as G28 U0 W0), feed, arcs, the thread cycle, the drill cycles, a dwell, stops, the tool change, the spindle, coolant, and more. A page has one \emph{template}. Pages that the post does not have of its own show a note: the post uses the default post for this. Tick the box to make it the post's own; it starts as a copy.
\end{description}

\subsection{Writing a template}
A template is the line or lines that the post writes for this thing. It is a text with \emph{tokens} between angle brackets. A token is replaced by a value: a code, a coordinate, a feed. For example:
\begin{center}\texttt{[<comp code> ]<code feed changed>[ X<X changed>][ Z<Z changed>][ F<F changed>]}\end{center}
is the template of a feed move in the LinuxCNC post. It writes \texttt{G1 X20 Z-5 F0.15}, and leaves out each value that did not change. What the pieces mean:
\begin{itemize}
  \item \texttt{<code feed>} is the code with the name \emph{feed} from the code table; \texttt{<X>} is the X coordinate. A token can have modifiers; \texttt{changed} writes it only when it differs from the line before.
  \item Square brackets make a part \emph{optional}. If the token inside has nothing to write, the whole part (including the space before the X) is left out.
\end{itemize}
Below the template, the page shows the same template written twice: with every optional part present and everything changed, and with nothing changed and every optional part gone. Look at both. A problem in the template is written in red.

Under the template is the list of tokens that work on this page. Double click a token, or click it, to put it into the template. \button{All tokens...} lists all of them.

\section{Making your own post}\label{sec:posts:own}
\begin{enumerate}
  \item In the editor, choose the post that is nearest to your controller and press \button{Save as copy...}. Give it a name.
  \item Go through the \emph{Settings} and the \emph{Code table}. Most differences between controllers are here.
  \item Go through the concept pages that matter: the tool change, the spindle and the cycles are the usual ones. Tick \emph{own} on a page only when you need to change it.
  \item Press \button{Save}. In \menupath{File}{Machines...}, choose the post in the machine.
\end{enumerate}
A post of your own is kept in the \file{posts} folder of the user folder (Chapter~\ref{ch:app_files}). You can copy the file to another PC.

\section{Checking what a post writes}\label{sec:posts:check}
Do this every time you make or change a post, before the machine runs it:
\begin{enumerate}
  \item Export the G-code of a small project that uses the things you changed (Section~\ref{sec:file:export}).
  \item Open the file in a text editor. Read it from the top. Check the first lines (units, plane, absolute or incremental), the tool change, the spindle start, one rapid move and one feed move.
  \item Run it on the machine without the part and without the tool, or at a safe height, with your hand on the stop.
\end{enumerate}

\chapter{Checking the program}\label{ch:checking}

WOP checks the program before the machine does. It looks for collisions, estimates the power and compares the speeds with the machine. It also lets you watch the program move by move. Nothing replaces your own look at the first part, but these checks find most mistakes while the part is still a picture.

\section{The rule: written as drawn}
WOP plans what you ask for. It does not shorten a cut because a tool might not reach, and it does not skip a path because it would collide. Instead, the check reports it. So a collision in the messages is not an error of WOP: it is a true thing about your program, and the G-code would do it.

Roughing is the one exception: it keeps the body of the tool off the finished part, and a line the tool cannot cut all the way is cut as far as it reaches. A note says what stays (Chapter~\ref{ch:blocks}). Finishing always follows the contour. If your tool cannot reach a place of the contour, the check reports a collision, and you turn the finishing pass of that block off.

\section{Collision check}\label{sec:checking:collision}
The check follows every move of the program. A cutting move removes the material the tool sweeps; other parts of the tool must stay out of the material. It reports these findings:
\begin{center}
\begin{tabularx}{\textwidth}{@{}l Y@{}}
\toprule
\textbf{Message} & \textbf{What to do} \\
\midrule
cuts into the finished part & Use a tool with a narrower body or more clearance, or change the block. \\
holder runs into the material & Use a tool with a smaller holder, or shorten the cut. \\
touches the jaws & Shorten the cut, or let the part stick out further. \\
rapid move through material & Check the safe positions and the approach of the block. \\
goes behind the clamping plane of the Setup & The part is held there. Shorten the cut, or change the clamping position in the Setup. \\
cutting move with the spindle stopped & Start the spindle before the block. A \emph{spindle stop} machine function may be in the way. \\
\bottomrule
\end{tabularx}
\end{center}
Touching is not a collision: only a deeper overlap is reported. A tool whose body reaches below the centre line is also tested on the far side of the part; the picture draws the mirrored material faintly.

A collision shows in three places: the block in the flow list gets a warning mark, the messages box has a line that starts with \textsf{COLLISION:}, and the picture shows a red cross where it happens. The line names the block and the move, with the same numbers that the step bar shows, and then the radius and Z. Step through the block: the cross appears with the move that causes it. From that block on, the picture shows the part with the material missing that the tool cut away.

\section{Machine checks}
These always run. They compare the program with the machine you chose in the Setup:
\begin{itemize}
  \item the spindle speed against the minimum and maximum RPM;
  \item the feed of each move against the fastest feed of the X and Z axis;
  \item the size of the part and of the cuts against the travel of the axes;
  \item whether the machine can do what a block needs: constant surface speed, or spindle-synchronised moves for threading.
\end{itemize}
Values that are not set in the machine are not checked.

\section{Power estimate}\label{sec:checking:power}
For each block WOP works out the power the cut needs: the cross section of the chip, times the feed, times the cutting speed at the largest radius that the block cuts, times the specific cutting force of the material (Chapter~\ref{ch:materials_menu}), with an efficiency of 80\,\%. If the number is higher than the power of the spindle of the machine, the messages say so, for example \emph{needs about 3.2 kW, the spindle has 2.2 kW; reduce the depth of cut or the feed}.

This is an estimate. It is meant to catch the program that is clearly too much, not to give you the exact power. For a block without a depth of cut (a custom path), it assumes twice the mean thickness of the material that the path takes off. If the material of the raw part is not in the table, or the machine has no spindle power, the power check says it was skipped.

\section{Step through}\label{sec:checking:step}
Step through (Section~\ref{sec:window:stepping}) is the best check there is. The tool is drawn at every move, with the material that is left at that moment. Go through a block you have not used before, and look at where the tool is at the places where you are not sure: the first approach, the corner of a groove, the end of a thread. If the tool collides in a move, the text next to the slider says so.

\section{The 3D view}\label{sec:checking:3d}
The 3D view (Section~\ref{sec:window:picture}) shows the part as a solid with a quarter cut away, so you see the cut face, the chuck, the jaws, the tool and the path. Drag with the mouse to turn it, use the wheel to zoom. It helps when the 2D section does not show you what you want: a groove at the far side, or the way the jaws hold the part.

\section{Reading the messages}\label{sec:checking:messages}
The messages box starts each line with a word (Section~\ref{sec:window:messages}). Work from top to bottom:
\begin{enumerate}
  \item \textsf{PROBLEM:} first. The program cannot be planned at that block. Fix the value it names.
  \item \textsf{COLLISION:} next. Each line names the block, the move and a way out.
  \item \textsf{note:} last. They are for your attention: a place that the tool did not reach, a value that was changed to fit.
\end{enumerate}
The checks run by themselves after every change, in the background: the picture comes first and the findings follow. While they run, the messages box says \textsf{WORKING:}.

\begin{tipbox}
A program with a collision may still be what you want, for example a finishing pass of an undercut that you cut by hand afterwards. WOP asks you at the export and lets you go on. But the machine will do exactly what the picture shows.
\end{tipbox}

\chapter{Exporting G-code}\label{ch:exporting}

The export itself is described in the chapter on the File menu (Section~\ref{sec:file:export}). This chapter is about what to do before it and after it.

\section{Before you export}
Go through this list once for every new program.
\begin{enumerate}
  \item \textbf{The Setup.} Is the machine the right one? Is the raw part the size of the bar you really have? Are the safe position and the tool change position safe on your lathe (Section~\ref{sec:project:safe})?
  \item \textbf{The tools.} Do the tool numbers match the turret? Is the tool offset right (Section~\ref{sec:tools:view})?
  \item \textbf{The checks.} They run by themselves after every change; wait until \textsf{WORKING:} is gone. Read all messages (Section~\ref{sec:checking:messages}). A line you accept is fine; a line you do not understand is not.
  \item \textbf{Step through.} Go through every block you have not used before (Section~\ref{sec:checking:step}).
  \item \textbf{The postprocessor.} If it is new or changed, check what it writes (Section~\ref{sec:posts:check}).
  \item \textbf{Save the project.} Then export with \menupath{File}{Export G-code...} (\key{Ctrl}+\key{E}). If the check found something, WOP lists it and asks; the answer is yours.
\end{enumerate}

\begin{notebox}
\menupath{File}{Export selected blocks...} writes only the blocks you ticked. Use it to run one block at a time while you try out a new program, or to continue after a stop (Section~\ref{sec:file:selected}).
\end{notebox}

\section{Running the program on the machine}
How the file gets to the machine depends on your set-up: a USB stick, the network, or a shared folder. The file ends in the extension that the postprocessor sets (\file{.ngc} for LinuxCNC).

For the first run of a new program:
\begin{enumerate}
  \item Open the G-code in a text editor and read the first lines and the first tool change.
  \item Set the work offset of the machine to match the Setup: X zero is the centre line, and Z zero is the place that the Setup calls \field{Front z}, usually the front face of the part (Section~\ref{sec:project:holding}). The \field{Work offset} of the Setup must be the one you set on the lathe (Section~\ref{sec:project:settings}).
  \item Touch off the tools. The tool offset numbers must match those in the program.
  \item Run the program with the part removed or with the tool well away from it, with your hand on the stop. Use the single block mode and a low rapid override.
  \item Then cut the first part. Look again at the first cuts.
\end{enumerate}

\begin{warningbox}
WOP and its postprocessors were not tried on every machine or controller. Anything in this manual that says it was not confirmed on a real machine is exactly that. You are the one who knows what your lathe will do. Stay at the machine.
\end{warningbox}


\appendix
\part{Reference}
\chapter{Keyboard shortcuts}\label{ch:shortcuts}

\section{In the main window}
\begin{center}
\begin{tabularx}{0.9\textwidth}{@{}l Y@{}}
\toprule
\textbf{Keys} & \textbf{What happens} \\
\midrule
\key{Ctrl}+\key{N} & New project \\
\key{Ctrl}+\key{O} & Open a project \\
\key{Ctrl}+\key{S} & Save the project \\
\key{Ctrl}+\key{Shift}+\key{S} & Save the project under a new name \\
\key{Ctrl}+\key{E} & Export the G-code \\
\key{Ctrl}+\key{Q} & Quit \\
\key{Ctrl}+\key{Z} & Undo \\
\key{Ctrl}+\key{Y} or \key{Ctrl}+\key{Shift}+\key{Z} & Redo \\
\key{Ctrl}+\key{3} & Switch between the 2D picture and the 3D view \\
\key{Left}, \key{Right} & With \field{Step through} on: one move back or forward \\
\key{Home}, \key{End} & With \field{Step through} on: the start or the end of the chosen block \\
\bottomrule
\end{tabularx}
\end{center}

\section{In the sketch editor}
\begin{center}
\begin{tabularx}{0.9\textwidth}{@{}l Y@{}}
\toprule
\textbf{Keys or mouse} & \textbf{What happens} \\
\midrule
Left click & Selects, or places what the chosen tool draws. \\
\key{Ctrl} + left click & Adds to the selection, or takes from it. \\
Drag on empty space & Selects everything inside the box you draw. \\
Drag a selected thing & Moves it; the relations decide how far it may go. \\
Mouse wheel & Zooms around the mouse. \\
Middle button, drag & Moves the view. \\
Right click & Ends what you are drawing (for example a line), the tool stays on. \\
\key{Esc} & Ends what you are drawing and goes back to selecting. With no tool on, it leaves the stepping through a path. It never closes the window. \\
\key{Delete} & Deletes the selected pieces, or the chosen relation. \\
Double click on a number & Changes the value of a dimension. \\
\key{Ctrl}+\key{Z} & Undo \\
\key{Ctrl}+\key{Y} & Redo \\
\key{Left}, \key{Right} & In the sketch of a custom path: step back or forward through the pieces of the path \\
\bottomrule
\end{tabularx}
\end{center}

\chapter{Files and folders}\label{ch:app_files}

\section{Project files}
A project is one file that ends in \file{.wop.json}. It is a text file, but not made to be edited by hand. It holds the Setup, the raw part, the blocks, the contours, the tools and workholding of the project, and a copy of the machine. You can copy it, mail it, and open it on another PC.

\section{The user folder}
Everything that belongs to you and not to one project is in the user folder (Chapter~\ref{ch:installing}): the folder \file{user_data} next to \file{WOP.exe} if there is one, otherwise \texttt{\%APPDATA\%\textbackslash WOP}.

\begin{center}
\begin{tabularx}{\textwidth}{@{}l Y@{}}
\toprule
\textbf{File or folder} & \textbf{What it holds} \\
\midrule
\file{settings.json} & The settings of the window (size, panels, the last folders, the legend) and which library of each kind is in use. \\
\file{conventions.json} & The values of \menupath{File}{Preferences...}. \\
\file{machines} & One file for each machine, \file{name.machine.json}. \\
\file{libraries} & The libraries, in one folder for each kind (\file{tools}, \file{workholding}, \file{materials}). Each library is one file, \file{name.json}. \\
\file{posts} & Your own postprocessors, \file{name.post.ini}. \\
\file{crash.log}, \file{wop.log} & Messages for the developer, see below. \\
\bottomrule
\end{tabularx}
\end{center}

\section{Libraries}
The library \emph{Example} is part of WOP and has no file; it is always there. Every other library is a file in the folder of its kind. You can copy a library file to another PC and use \menupath{Tools}{Library} $\rightarrow$ \menu{Import library file...} to add it. \menu{Open library folder} in the same submenu shows the folder.

\section{Settings and preferences}
\file{settings.json} is written by WOP when you change something. \file{conventions.json} holds the Preferences. If you delete either file, WOP starts with its built-in values. Your machines and libraries stay.

\section{Postprocessor files}
The posts that come with WOP are in the program folder and cannot be changed. Your own posts are in the folder \file{posts} of the user folder. A post file is plain text (Chapter~\ref{ch:posts}).

\section{Log files}
\file{wop.log} holds the messages that Qt, the toolkit of the window, wrote. \file{crash.log} is written when WOP stops without a message, or when the window stopped answering for about two seconds. It holds what each part of the program was doing. If WOP closes by itself, keep \file{crash.log} and \file{wop.log} and give them to the developer.

\chapter{Glossary}\label{ch:glossary}

\begin{description}
\item[Block] One step of the program, for example facing or drilling. Blocks are done from top to bottom.
\item[Boring bar] A tool for work inside a bore.
\item[Chuck body] The body of the chuck that holds the jaws.
\item[Clamping plane] The plane in Z where the jaws hold the part. A cut behind it is reported.
\item[Clearance] The distance the tool keeps from the material while it moves inside a block.
\item[Collision] The tool, its holder or the jaws touching material or the part where they must not. WOP reports it; it does not stop the program.
\item[Constant surface speed] The spindle speed changes with the radius so that the cutting speed stays the same (G96 on LinuxCNC).
\item[Construction line] A line in the sketch editor that helps drawing but is not part of the shape.
\item[Contour] A shape drawn in the sketch editor from lines and arcs.
\item[Custom path] A block whose path you draw yourself in the sketch editor, with rapid and feed lines.
\item[Cutting edge] The part of a tool that is allowed to touch the material. It is drawn in orange.
\item[Dimension] A measure in the sketch editor, for example the length of a line. The drawing follows it.
\item[Flow list] The list of blocks on the left of the window.
\item[Fully defined] A sketch in which every point is fixed by relations and dimensions.
\item[Jaws] The pieces of the chuck that grip the part.
\item[Library] A file with tools, workholding or materials that you use in many projects.
\item[Locked] A thing in the sketch editor that you cannot change there because it comes from elsewhere (the safe point, the outline of the part).
\item[Machine] What belongs to your lathe: limits, speeds, power, tool change and postprocessor.
\item[Named point] A point with a name that WOP uses, for example the tip of a tool.
\item[Offset (tool)] The number the postprocessor calls after a tool change (G43 H...).
\item[Postprocessor] A text file that says how WOP writes G-code for a controller.
\item[Preferences] Settings that apply to every new project.
\item[Project] One job in one file.
\item[Radius and diameter] X is the distance from the centre line (radius). The forms can show it as a diameter (twice the radius).
\item[Raw part] The bar or tube you start with.
\item[Relation] A rule in the sketch editor, for example that two lines are parallel.
\item[Roughing] Cutting away the bulk of the material, with some material left for the finishing pass.
\item[Safe position] Where the tool waits between blocks, outside the part and the chuck.
\item[Setup] The first entry of the flow list: the machine, the raw part, the safe distances and the workholding of the project.
\item[Sketch editor] The window where shapes are drawn.
\item[Specific cutting force] A number of the material that, with the chip size, gives the cutting force. It is used for the power estimate.
\item[Stick-out] How far the part sticks out of the jaws.
\item[Style] How a block orders its cuts, for example the grooving styles standard and layers.
\item[Tool change position] Where the tool is changed.
\item[Toolbox] A set of library tools with their numbers in the turret.
\item[Touch off] Setting the zero of a tool on the machine.
\item[Work offset] The zero point of the part on the machine (G54 and the following).
\end{description}

\chapter{When something goes wrong}\label{ch:troubleshooting}

\section{The program does not start}
\begin{itemize}
  \item Start the \file{WOP.exe} that is in the program folder. A copy outside the folder cannot find its parts.
  \item If a virus scanner stops it, allow the folder.
  \item If the window opens and closes at once, look for \file{crash.log} and \file{wop.log} in the user folder (Chapter~\ref{ch:app_files}).
\end{itemize}

\section{WOP closes by itself}
If WOP closes with no message, the program stopped. Keep \file{crash.log} and \file{wop.log} from the user folder and tell the developer what you did just before. A note of the steps is the most useful thing. WOP writes the two files itself; you do not need to turn anything on.

\section{The 3D button is grey}
Your graphics card or its driver does not give OpenGL 2. The hover text of the button says why. Update the graphics driver. Everything else works in 2D.

\section{A message in the flow list}
\begin{description}
\item[A warning mark at a block] The check found a collision, or the block has a problem. Select the block and read the messages box (Section~\ref{sec:checking:messages}).
\item[\textsf{PROBLEM:}] The block cannot be planned. The text says which value.
\item[\textsf{WORKING:}] Wait. The picture is that of the program before your last change until the work is done.
\item[Machine without CSS or synchronised moves] The block needs a function that is not ticked in the machine (Section~\ref{sec:machines:settings}). Tick it if your controller has it; otherwise change the block.
\item[Power] The cut needs more power than the spindle has. Reduce the depth of cut or the feed (Section~\ref{sec:checking:power}).
\end{description}

\section{A value is refused}
A field that does not accept your value gives a message and keeps the old value. Typing a diameter where WOP wants a radius is common: look at the label of the field and at the X mode of the Setup (Section~\ref{sec:project:xmode}).

\section{The sketch editor}
\begin{description}
\item[I cannot move a point] The sketch has a relation or a dimension that holds it, or the point is locked (Section~\ref{sec:sketch:locked}).
\item[The editor says there is a problem] The shape breaks a rule of its use, for example it is not closed. The text in the editor says what (Section~\ref{sec:sketch:window}).
\end{description}

\section{The G-code is not what I expected}
\begin{itemize}
  \item Is the right machine in the Setup, and does it use the right postprocessor (Chapter~\ref{ch:machines})?
  \item Is the post known to work for your controller? Read what it writes (Section~\ref{sec:posts:check}).
  \item Does the program show the right path in the picture? WOP writes what the picture shows.
  \item Is X a diameter in the G-code? That is a setting of the post (\field{x diameters}).
  \item Does the program start where you thought? Look at the safe position and the tool change position in the Setup (Section~\ref{sec:project:safe}).
\end{itemize}

\section{A tool or the library asks a question}
When you close or start a new project, WOP may ask whether to update the tool library with tools that you changed (Section~\ref{sec:file:questions}). \button{Leave the library as it is} keeps the library; the project keeps its own tools.


\end{document}
